\documentclass[11pt]{article}
\usepackage[T1]{fontenc}
\usepackage[utf8]{inputenc}
\usepackage{lmodern}
\usepackage[margin=1in]{geometry}
\usepackage{amsmath,amssymb,mathtools}
\usepackage{booktabs,longtable,array}
\usepackage{fancyvrb}
\usepackage[hidelinks]{hyperref}
\usepackage{microtype}
\pdfobjcompresslevel=0
\pdfinfoomitdate=1
\pdftrailerid{}
\pdfsuppressptexinfo=15
\setlength{\parindent}{0pt}
\setlength{\parskip}{0.65em}
\setlength{\emergencystretch}{3em}
\title{Pairing of universes of masses +m and -m: exact proofs for dirac16complex and dirac16complex00\\[0.5em]\large The pairing theorems T1 and T2 and the quantum-level reading in the author's primordial gravitational field, with every hypothesis, the verification records and an exact statement of what they do not establish}
\author{dirac16complex project (with Claude Opus 5.5)}
\date{September 2026}
\begin{document}
\maketitle
\tableofcontents
\newpage
\begin{abstract}
The author asked on 2026-10-01: ``PROVE that Universes of masses \{+mass, -mass\} are created in pairs for each case of the dirac16complex and the dirac16complex00 fields.'' This document gives the exact answer computed by the code in \texttt{Revision\allowbreak{}/} for the author's primordial metric. Proved, for both fields and with every hypothesis stated: T1 (chirality pairing): the chirality matrix $\Gamma$ maps every configuration $\Psi$ of the theory with mass and coupling $(m,\lambda)$ to the configuration $\Gamma\Psi$ of the theory with $(-m,-\lambda)$, with $\mathcal{L}_{m,\lambda}[\Gamma\Psi] = -\mathcal{L}_{-m,-\lambda}[\Psi]$, solutions to solutions and the energy-momentum tensor and the current reversed, so that the pair has zero total energy-momentum, current and charge as classical bilinears, in every gravitational field. T2 (mirror pairing): $\Gamma$ combined with a Pin(4,4) reflection of character $-1$ maps $(m,\lambda)$ to $(-m,\lambda)$ with $\mathcal{L} \to +\mathcal{L}$; in the author's field it is the mirror across the Z2 brane $z = \pi/2$ (an ASSUMED construction), and the partner has EQUAL, not opposite, energy-momentum. The quantum-level reading for dirac16complex: the chirality image carries the Krein metric $-B$ and is the same quantum system re-labelled; two independently quantised universes of masses $+m$ and $-m$ have identical one-particle spectra and their generators do not cancel. A corollary: a T1 pair taken as the complete classical source of the author's metric is a zero source, and the Einstein equations then have no solution for $H > 0$. The Kohn-Sham level T3 is not established here (it is to be added; its verified ingredients are pointed to). Not established by any of these equations: a creation process, a rate, a probability or an amplitude for creating universes, in pairs or otherwise. The pairing theorems are proved; the creation of pairs is not.

\end{abstract}

\section{The request and the exact answer}

The author's request, quoted from \texttt{Revision\allowbreak{}/README.\allowbreak{}md}: ``PROVE that Universes of masses \{+mass, -mass\} are created in pairs for each case of the dirac16complex and the dirac16complex00 fields.''

What a universe of mass $m$ means in this record: a configuration of one of the two fields with mass parameter $m$ and coupling $\lambda$ in the author's primordial gravitational field, i.e. a solution of the Euler-Lagrange equations of section 2.4. Two such universes are paired when an explicit invertible map takes every solution of the theory with mass $m$ to a solution of the theory with mass $-m$ and takes its Lagrangian, energy-momentum tensor and current to stated multiples of those of the partner.

The exact answer, for both fields:

\begin{enumerate}
\item Proved (sections 4 to 7): the pairing theorems T1 and T2, the quantum-level reading Q for dirac16complex and the corollary C1 for the field equations of $a_4$, each with its exact hypotheses.
\item Not proved, and not derivable from these equations: that such universes are CREATED. The equations are field equations on a fixed gravitational background and their canonical quantisation; no equation of this record produces a pair of universes from anything, and no creation process, rate or amplitude follows from them (section 11).
\item The Kohn-Sham level T3 of \texttt{Revision\allowbreak{}/SPEC.md} section 9 is not established in this document; section 8 states it and points to the Kohn-Sham ingredients that are already verified.
\end{enumerate}

Every formula and number below is taken from the Revision reports listed in section 9; nothing is taken from the earlier stages of the repository. Interpretations are labelled as such.

\section{Setting and conventions}

\subsection{The author's metric and coordinates}

The coordinates are named as the author names them: $x_1, x_2, x_3$ are ordinary 3-space (inflating), $x_4$ is the time, $x_5, x_6, x_7$ are the three extra times, which DEFLATE exponentially as $a_4$ increases, and $x_8$ is the hidden space direction, with $z = 6Hx_8 \in (0,\pi/2)$ and $H > 0$ the author's constant. The metric of \texttt{Revision\allowbreak{}/SPEC.md} section 1 is

\[
\begin{aligned}
ds^2 = {} & e^{2a_4(x_4)}\sin^{1/3}z\,\big(dx_1^2+dx_2^2+dx_3^2\big) - dx_4^2 \\
& - e^{-2a_4(x_4)}\sin^{1/3}z\,\big(dx_5^2+dx_6^2+dx_7^2\big) + \cot^2 z\,dx_8^2 ,
\end{aligned}
\]

of signature (4,4) (space-like $x_1, x_2, x_3, x_8$; time-like $x_4, x_5, x_6, x_7$), with $\sqrt{|g|} = \cos z$. The diagonal vielbein is $e^a{}_\mu = f_a\delta^a_\mu$ with $f_1 = f_2 = f_3 = e^{a_4}\sin^{1/6}z$, $f_4 = 1$, $f_5 = f_6 = f_7 = e^{-a_4}\sin^{1/6}z$ and $f_8 = \cot z$; the frame metric is $\eta = \mathrm{diag}(+1,+1,+1,-1,-1,-1,-1,+1)$ in the order $x_1, \dots, x_8$. Frame indices $a = 1, \dots, 8$ are aligned with the coordinates, so $\gamma^4$ is the time gamma $\gamma^{(x_4)}$ and $\gamma^8$ the hidden-direction gamma $\gamma^{(x_8)}$.

The canonical spin connection (vielbein postulate) has the nonzero components $\omega_{x_i,\,i4} = a_4' e^{a_4}\sin^{1/6}z$ and $\omega_{x_i,\,i8} = H e^{a_4}\sin^{1/6}z$ for the inflating directions $i = 1, 2, 3$, and $\omega_{x_t,\,4t} = -a_4' e^{-a_4}\sin^{1/6}z$ and $\omega_{x_t,\,t8} = -H e^{-a_4}\sin^{1/6}z$ for the deflating extra times $t = 5, 6, 7$, with $\Omega_{x_4} = \Omega_{x_8} = 0$. Contracted with the gammas, $\gamma^\mu\Omega_\mu = 3H\gamma^8$: each inflating direction contributes $\frac{a_4'}{2}\gamma^4 + \frac{H}{2}\gamma^8$ and each deflating extra time $-\frac{a_4'}{2}\gamma^4 + \frac{H}{2}\gamma^8$, so the time-direction terms cancel and the hidden-direction terms add.

\subsection{Clifford data}

The gammas are the author's real $16 \times 16$ matrices T16, re-constructed in Revision code from the author's tau formulas, with $\gamma^8 = T16[0]$, $\gamma^{1,2,3} = T16[1..3]$, $\gamma^4 = T16[4]$ and $\gamma^{5,6,7} = T16[5..7]$. They are real signed permutation matrices with $\{\gamma^a,\gamma^b\} = 2\eta^{ab}$ and $(\gamma^a)^T = \eta_{aa}\gamma^a$. Further:

\begin{itemize}
\item $C = \gamma^8\gamma^1\gamma^2\gamma^3$ is real and symmetric, $C^2 = 1$, $C\gamma^a$ is real and antisymmetric and $C\gamma^aC^{-1} = -(\gamma^a)^T$; the Dirac adjoint is $\bar\Psi = \Psi^\dagger C$.
\item $\Gamma = \gamma^8\gamma^1\gamma^2\gamma^3\gamma^4\gamma^5\gamma^6\gamma^7 = \mathrm{diag}(-I_8, I_8)$.
\item $B = -iC\gamma^4$ is Hermitian with $B^2 = 1$ and signature (8,8).
\item $S^{ab} = \frac14[\gamma^a,\gamma^b]$ satisfy the commutation relations of so(4,4).
\item The 16-dimensional representation is irreducible under Pin(4,4) (commutant of dimension 1) and splits under Spin(4,4) into two inequivalent irreducible 8-dimensional representations, the chiral halves $\Gamma = -1$ (components 1 to 8, written $\psi_-$) and $\Gamma = +1$ (components 9 to 16, written $\psi_+$); every $\gamma^a$ exchanges the two halves.
\end{itemize}

\subsection{The two fields and the Lagrangian}

dirac16complex is $\Psi$: 16 complex ANTICOMMUTING (Grassmann) components, a Pin(4,4) spinor, quantised canonically (section 6). dirac16complex00 is $\Phi$: 16 complex COMMUTING components that transform as a Pin(4,4) spinor, a classical field. Both are coupled to gravity through the vielbein and the canonical spin connection. The theorems are stated with the letter $\Psi$ for both fields; ``both statistics'' means that the statement and its proof hold for Grassmann components and for commuting components. With $S = \bar\Psi\Psi$, $D_\mu\Psi = \partial_\mu\Psi + \Omega_\mu\Psi$, $D_\mu\bar\Psi = \partial_\mu\bar\Psi - \bar\Psi\Omega_\mu$, $\Omega_\mu = \frac12\omega_{\mu ab}S^{ab}$ and $\gamma^\mu = e^\mu{}_a\gamma^a$, the Lagrangian density of both fields is

\[
\mathcal{L}_{m,U}[\Psi] = \sqrt{|g|}\,\big[K[\Psi] - mS - U(S)\big],\qquad
K[\Psi] = \tfrac12\big(\bar\Psi\gamma^\mu D_\mu\Psi - (D_\mu\bar\Psi)\gamma^\mu\Psi\big),
\]

written $\mathcal{L}_{m,\lambda}$ for $U(S) = \frac{\lambda}{2}S^2$. It is real for both statistics.

\subsection{Field equation, energy-momentum tensor, current}

The Euler-Lagrange equations of $\mathcal{L}_{m,U}$ are, for both statistics,

\[
E_{m,U}[\Psi] \equiv \gamma^\mu D_\mu\Psi - \big(m + U'(S)\big)\Psi = 0,\qquad
(D_\mu\bar\Psi)\gamma^\mu = -\big(m + U'(S)\big)\bar\Psi .
\]

The symmetric (Belinfante) energy-momentum tensor of the theory record is

\[
T^\nu{}_\mu = \delta^\nu_\mu\,\frac{\mathcal{L}}{\sqrt{|g|}} - \frac14\Big(\bar\Psi\gamma^\nu D_\mu\Psi - (D_\mu\bar\Psi)\gamma^\nu\Psi + \bar\Psi\gamma_\mu D^\nu\Psi - (D^\nu\bar\Psi)\gamma_\mu\Psi\Big),
\]

with the energy density $\rho = -T^{x_4}{}_{x_4}$ and the pressures $p_\mu = T^\mu{}_\mu$ (no sum) of \texttt{Revision\allowbreak{}/SPEC.md} section 4; for homogeneous on-shell states $\rho = mS + U$ and $p = SU' - U$. The conserved current is $J^\mu = -i\bar\Psi\gamma^\mu\Psi$, with the charge density $J^{x_4} = \Psi^\dagger B\Psi$ and the charge $Q = \int\cos z\,\Psi^\dagger B\Psi\,d^7x$.

Convention note. The pairing records use the tensor

\[
T_{\mu\nu} = \tfrac14\big(\bar\Psi\gamma_\mu D_\nu\Psi + \bar\Psi\gamma_\nu D_\mu\Psi - (D_\mu\bar\Psi)\gamma_\nu\Psi - (D_\nu\bar\Psi)\gamma_\mu\Psi\big) - g_{\mu\nu}\,\mathcal{L}/\sqrt{|g|},
\]

which is minus the tensor above, and the currents $\bar\Psi\gamma^\mu\Psi$ (Wolfram) and $i\bar\Psi\gamma^\mu\Psi$ (sympy). Every statement of T1, T2 and Q about $T$ and $J$ is linear and homogeneous in them ($T' = \pm T$, $J' = \pm J$), so it holds unchanged for every overall sign and constant factor.

\subsection{The field equation in the author's metric}

Written out in the author's metric, the field equation of both fields is

\[
\begin{aligned}
& e^{-a_4}\sin^{-1/6}z\,\big(\gamma^1\partial_1+\gamma^2\partial_2+\gamma^3\partial_3\big)\Psi + \gamma^4\partial_4\Psi
+ e^{a_4}\sin^{-1/6}z\,\big(\gamma^5\partial_5+\gamma^6\partial_6+\gamma^7\partial_7\big)\Psi \\
& \qquad + \tan z\,\gamma^8\partial_8\Psi + 3H\gamma^8\Psi = \big(m + U'(S)\big)\Psi ,
\end{aligned}
\]

where $3H\gamma^8 = \gamma^\mu\Omega_\mu$ is the spin-connection term; it is nonzero for every $H > 0$ and every $\Psi \neq 0$ (non-triviality [1] and [2] of the theory record). In the chiral block form $\Psi = (\psi_-, \psi_+)$ every gamma is block off-diagonal, with upper-right block $\bar\tau^a$ and lower-left block $\tau^a$, and $\bar\tau^8 = \tau^8 = I_8$; the 16 equations become the two 8-component equations

\[
\sum_a f_a^{-1}\bar\tau^a\partial_a\psi_+ + 3H\psi_+ = V\psi_-,\qquad
\sum_a f_a^{-1}\tau^a\partial_a\psi_- + 3H\psi_- = V\psi_+,\qquad V = m + U'(S).
\]

\subsection{Records for section 2}

\begingroup
\small
\begin{longtable}{@{}>{\raggedright\arraybackslash}p{0.267\linewidth}>{\raggedright\arraybackslash}p{0.267\linewidth}>{\raggedright\arraybackslash}p{0.267\linewidth}@{}}
\toprule
\textbf{item} & \textbf{Wolfram records} & \textbf{sympy records} \\
\midrule
\endfirsthead
\toprule
\textbf{item} & \textbf{Wolfram records} & \textbf{sympy records} \\
\midrule
\endhead
metric, vielbein, $\sqrt{\lvert g\rvert} = \cos z$ & \texttt{metric\_i\allowbreak{}s\_the\_au\allowbreak{}thors}, \texttt{sqrt\_det\allowbreak{}\_g\_is\_co\allowbreak{}s\_z}, \texttt{primordi\allowbreak{}al\_vielb\allowbreak{}ein} & \texttt{metric\_f\allowbreak{}rom\_viel\allowbreak{}bein\_equ\allowbreak{}als\_SPEC}, \texttt{sqrt\_det\allowbreak{}\_g\_equal\allowbreak{}s\_cos\_z} \\
spin connection and $\gamma^\mu\Omega_\mu = 3H\gamma^8$ & \texttt{omega\_co\allowbreak{}mponents}, \texttt{Omega\_co\allowbreak{}mponents}, \texttt{gammaOme\allowbreak{}ga\_equal\allowbreak{}s\_3H\_gam\allowbreak{}ma\_x8}, \texttt{gammaOme\allowbreak{}ga\_x4\_te\allowbreak{}rms\_canc\allowbreak{}el} & \texttt{geometry\allowbreak{}.spin\_co\allowbreak{}nnection\allowbreak{}\_compone\allowbreak{}nts}, \texttt{gamma\_mu\allowbreak{}\_Omega\_m\allowbreak{}u\_equals\allowbreak{}\_3H\_gamm\allowbreak{}a\_x8}, \texttt{time\_ter\allowbreak{}ms\_cance\allowbreak{}l\_hidden\allowbreak{}\_term\_su\allowbreak{}rvives} \\
gammas & \texttt{coordina\allowbreak{}te\_map}, \texttt{Clifford\allowbreak{}\_relatio\allowbreak{}n}, \texttt{\detokenize{reality}}, \texttt{signed\_p\allowbreak{}ermutati\allowbreak{}on\_matri\allowbreak{}ces}, \texttt{symmetry\allowbreak{}\_pattern} & \texttt{clifford\allowbreak{}\_relatio\allowbreak{}n}, \texttt{reality\_\allowbreak{}signed\_p\allowbreak{}ermutati\allowbreak{}ons}, \texttt{symmetry\allowbreak{}\_pattern} \\
$C$, $\Gamma$, $B$, $S^{ab}$ & \texttt{C\_real\_s\allowbreak{}ymmetric}, \texttt{C\_square\allowbreak{}d\_identi\allowbreak{}ty}, \texttt{C\_gamma\_\allowbreak{}real\_ant\allowbreak{}isymmetr\allowbreak{}ic}, \texttt{C\_gamma\_\allowbreak{}C\_invers\allowbreak{}e}, \texttt{Gamma\_de\allowbreak{}finition}, \texttt{Gamma\_di\allowbreak{}ag}, \texttt{B\_Hermit\allowbreak{}ian}, \texttt{B\_square\allowbreak{}d\_identi\allowbreak{}ty}, \texttt{B\_signat\allowbreak{}ure\_8\_8}, \texttt{S\_Lorent\allowbreak{}z\_algebr\allowbreak{}a} & \texttt{C\_real\_s\allowbreak{}ymmetric\allowbreak{}\_involut\allowbreak{}ion}, \texttt{C\_gamma\_\allowbreak{}antisymm\allowbreak{}etric}, \texttt{C\_conjug\allowbreak{}ation}, \texttt{chiralit\allowbreak{}y\_diag}, \texttt{B\_hermit\allowbreak{}ian\_invo\allowbreak{}lution\_s\allowbreak{}ignature}, \texttt{S\_lorent\allowbreak{}z\_algebr\allowbreak{}a} \\
Pin(4,4) and Spin(4,4) & \texttt{Pin44\_ir\allowbreak{}reducibl\allowbreak{}e\_commut\allowbreak{}ant\_dim\_\allowbreak{}1}, \texttt{Spin44\_c\allowbreak{}ommutant\allowbreak{}\_dim\_2\_c\allowbreak{}hiral\_pr\allowbreak{}ojectors}, \texttt{chiral\_h\allowbreak{}alves\_ir\allowbreak{}reducibl\allowbreak{}e}, \texttt{chiral\_h\allowbreak{}alves\_in\allowbreak{}equivale\allowbreak{}nt\_inter\allowbreak{}twiners\_\allowbreak{}0}, \texttt{Spin\_Pin\allowbreak{}\_reflect\allowbreak{}ion\_swap\allowbreak{}s\_halves} & \texttt{pin\_comm\allowbreak{}utant\_di\allowbreak{}mension\_\allowbreak{}1}, \texttt{spin\_com\allowbreak{}mutant\_d\allowbreak{}imension\allowbreak{}\_2}, \texttt{spin\_hal\allowbreak{}ves\_irre\allowbreak{}ducible}, \texttt{spin\_hal\allowbreak{}ves\_ineq\allowbreak{}uivalent}, \texttt{reflecti\allowbreak{}ons\_exch\allowbreak{}ange\_hal\allowbreak{}ves} \\
reality of $\mathcal{L}$ & \texttt{\detokenize{L_real_G}}, \texttt{\detokenize{L_real_C}} & \texttt{grassman\allowbreak{}n\_lagran\allowbreak{}gian\_rea\allowbreak{}l}, \texttt{commutin\allowbreak{}g\_lagran\allowbreak{}gian\_rea\allowbreak{}l} \\
Euler-Lagrange equations & \texttt{\detokenize{EL_Psi_G}}, \texttt{EL\_Psiba\allowbreak{}r\_G}, \texttt{\detokenize{EL_Psi_C}}, \texttt{EL\_Psiba\allowbreak{}r\_C}, \texttt{T1\_Euler\allowbreak{}\_Lagrang\allowbreak{}e\_derive\allowbreak{}d\_primor\allowbreak{}dial\_gra\allowbreak{}ssmann}, \texttt{T1\_Euler\allowbreak{}\_Lagrang\allowbreak{}e\_derive\allowbreak{}d\_primor\allowbreak{}dial\_com\allowbreak{}muting} & \texttt{grassman\allowbreak{}n\_euler\_\allowbreak{}lagrange\allowbreak{}\_psi\_var\allowbreak{}iation}, \texttt{grassman\allowbreak{}n\_euler\_\allowbreak{}lagrange\allowbreak{}\_psibar\_\allowbreak{}variatio\allowbreak{}n}, \texttt{commutin\allowbreak{}g\_euler\_\allowbreak{}lagrange\allowbreak{}\_psi\_var\allowbreak{}iation}, \texttt{commutin\allowbreak{}g\_euler\_\allowbreak{}lagrange\allowbreak{}\_psibar\_\allowbreak{}variatio\allowbreak{}n} \\
energy-momentum tensor & \texttt{T\_symmet\allowbreak{}ric\_part\allowbreak{}\_Belinfa\allowbreak{}nte\_G}, \texttt{T\_symmet\allowbreak{}ric\_part\allowbreak{}\_Belinfa\allowbreak{}nte\_C}, \texttt{kinetic\_\allowbreak{}sum\_on\_s\allowbreak{}hell\_G}, \texttt{kinetic\_\allowbreak{}sum\_on\_s\allowbreak{}hell\_C} & \texttt{grassman\allowbreak{}n\_emt\_sy\allowbreak{}mmetric}, \texttt{commutin\allowbreak{}g\_emt\_sy\allowbreak{}mmetric}, \texttt{grassman\allowbreak{}n\_homoge\allowbreak{}neous\_on\allowbreak{}\_shell\_r\allowbreak{}ho\_p}, \texttt{commutin\allowbreak{}g\_homoge\allowbreak{}neous\_on\allowbreak{}\_shell\_r\allowbreak{}ho\_p} \\
current and charge & \texttt{current\_\allowbreak{}conserva\allowbreak{}tion\_ide\allowbreak{}ntity\_G}, \texttt{current\_\allowbreak{}conserva\allowbreak{}tion\_ide\allowbreak{}ntity\_C}, \texttt{charge\_d\allowbreak{}ensity\_i\allowbreak{}s\_Krein\_\allowbreak{}form\_G}, \texttt{charge\_d\allowbreak{}ensity\_i\allowbreak{}s\_Krein\_\allowbreak{}form\_C} & \texttt{grassman\allowbreak{}n\_curren\allowbreak{}t\_conser\allowbreak{}vation}, \texttt{commutin\allowbreak{}g\_curren\allowbreak{}t\_conser\allowbreak{}vation} \\
field equation in the metric, block form, non-triviality & \texttt{Dirac\_op\allowbreak{}erator\_e\allowbreak{}xplicit\_\allowbreak{}G}, \texttt{Dirac\_op\allowbreak{}erator\_e\allowbreak{}xplicit\_\allowbreak{}C}, \texttt{block\_fo\allowbreak{}rm}, \texttt{nontrivi\allowbreak{}ality\_1\_\allowbreak{}dirac16c\allowbreak{}omplex}, \texttt{nontrivi\allowbreak{}ality\_2\_\allowbreak{}dirac16c\allowbreak{}omplex00} & \texttt{nontrivi\allowbreak{}ality\_Om\allowbreak{}ega\_zero\allowbreak{}\_iff\_fla\allowbreak{}t} \\
\bottomrule
\end{longtable}
\endgroup

\section{Algebraic lemmas}

All four lemmas are exact matrix identities of the Revision gammas, verified by WolframScript and independently by sympy (records after each proof).

\textbf{Lemma 1 (the chirality matrix).} $\Gamma$ is real and symmetric, $\Gamma^2 = 1$, $\Gamma\gamma^a = -\gamma^a\Gamma$ for every $a$, $\Gamma C = C\Gamma$ and $\Gamma S^{ab} = S^{ab}\Gamma$. Hence $\Gamma\Omega_\mu = \Omega_\mu\Gamma$ for EVERY connection $\Omega_\mu = \frac12\omega_{\mu ab}S^{ab}$, and $D_\mu(\Gamma\Psi) = \Gamma D_\mu\Psi$, $D_\mu(\bar\Psi\Gamma) = (D_\mu\bar\Psi)\Gamma$.

Proof. $\Gamma$ is the product of all eight different $\gamma^a$. Moving $\gamma^b$ through it passes the seven factors $\gamma^a$, $a \neq b$, which anticommute with $\gamma^b$, and the factor $\gamma^b$ itself, which commutes with it: the sign is $(-1)^7 = -1$. $C$ is a product of four and $S^{ab}$ ($a \neq b$) a multiple of a product of two different gammas, so both commute with $\Gamma$. $\Gamma = \mathrm{diag}(-I_8, I_8)$ is computed exactly; it is real, symmetric and squares to 1. $\Gamma$ is constant and commutes with $\Omega_\mu$, which gives the statements about $D_\mu$. QED.

\begingroup
\small
\begin{longtable}{@{}>{\raggedright\arraybackslash}p{0.400\linewidth}>{\raggedright\arraybackslash}p{0.400\linewidth}@{}}
\toprule
\textbf{Wolfram records} & \textbf{sympy records} \\
\midrule
\endfirsthead
\toprule
\textbf{Wolfram records} & \textbf{sympy records} \\
\midrule
\endhead
\texttt{Gamma\_pr\allowbreak{}operties}, \texttt{Gamma\_di\allowbreak{}ag}, \texttt{Gamma\_an\allowbreak{}ticommut\allowbreak{}es\_with\_\allowbreak{}gammas}, \texttt{Gamma\_co\allowbreak{}mmutes\_w\allowbreak{}ith\_C}, \texttt{S\_commut\allowbreak{}es\_with\_\allowbreak{}Gamma} & \texttt{gammas.G\allowbreak{}amma}, \texttt{gammas.G\allowbreak{}amma\_ant\allowbreak{}icommute\allowbreak{}s}, \texttt{chiralit\allowbreak{}y\_antico\allowbreak{}mmutes}, \texttt{chiralit\allowbreak{}y\_C\_rela\allowbreak{}tion}, \texttt{S\_preser\allowbreak{}ves\_C\_an\allowbreak{}d\_commut\allowbreak{}es\_with\_\allowbreak{}Gamma} \\
\bottomrule
\end{longtable}
\endgroup

\textbf{Lemma 2 (bilinears under the chirality map).} For $\Psi' = \Gamma\Psi$: $\bar\Psi' = \bar\Psi\Gamma$; $\bar\Psi' M\Psi' = (-1)^k\bar\Psi M\Psi$ for every product $M$ of $k$ gammas; in particular $S$ is invariant and every bilinear with $\gamma^a$, $\gamma^aS^{bc}$ or $S^{bc}\gamma^a$ changes sign. Moreover $\Gamma B\Gamma = -B$.

Proof. By Lemma 1, $\bar\Psi' = \Psi^\dagger\Gamma^\dagger C = \Psi^\dagger\Gamma C = \Psi^\dagger C\Gamma = \bar\Psi\Gamma$. Then $\bar\Psi'M\Psi' = \bar\Psi\,\Gamma M\Gamma\,\Psi$, and $\Gamma M\Gamma = (-1)^kM$ by Lemma 1. Finally $\Gamma B\Gamma = -iC\,\Gamma\gamma^4\Gamma = iC\gamma^4 = -B$. QED.

\begingroup
\small
\begin{longtable}{@{}>{\raggedright\arraybackslash}p{0.400\linewidth}>{\raggedright\arraybackslash}p{0.400\linewidth}@{}}
\toprule
\textbf{Wolfram records} & \textbf{sympy records} \\
\midrule
\endfirsthead
\toprule
\textbf{Wolfram records} & \textbf{sympy records} \\
\midrule
\endhead
\texttt{T1\_kerne\allowbreak{}l\_scalar}, \texttt{T1\_kerne\allowbreak{}l\_kineti\allowbreak{}c}, \texttt{T1\_kerne\allowbreak{}l\_connec\allowbreak{}tion} (all 512 triples), \texttt{T1\_kerne\allowbreak{}l\_field\_\allowbreak{}equation}, \texttt{Q\_Krein\_\allowbreak{}metric\_o\allowbreak{}f\_images} & \texttt{T1.gener\allowbreak{}al\_field\allowbreak{}.matrix\_\allowbreak{}identiti\allowbreak{}es} (456 kinetic and connection basis matrices), \texttt{Q.image\_\allowbreak{}krein\_me\allowbreak{}tric} \\
\bottomrule
\end{longtable}
\endgroup

\textbf{Lemma 3 (reflections and their Pin(4,4) lifts).} Let $n$ be a frame direction and $R_n$ the diagonal matrix with $-1$ in place $n$ and $+1$ elsewhere (the reflection of that direction; $R_n\eta R_n = \eta$).

\begin{enumerate}
\item $(\gamma^n)^TC\gamma^n = -\eta_{nn}C$.
\item $\gamma^n\gamma^a(\gamma^n)^{-1} = -(R_n)^a{}_b\gamma^b$ and $\gamma^nS^{ab}(\gamma^n)^{-1} = (R_n)^a{}_c(R_n)^b{}_dS^{cd}$.
\item $P_n := \Gamma\gamma^n$ is plus or minus the product of the seven other gammas, hence an element of Pin(4,4); $P_n^{-1}\gamma^aP_n = (R_n)^a{}_b\gamma^b$ (it covers $R_n$); $P_n^\dagger CP_n = \chi_nC$ with the character $\chi_n = -\eta_{nn}$, i.e. $-1$ for the space-like $x_1, x_2, x_3, x_8$ and $+1$ for the time-like $x_4, x_5, x_6, x_7$; and $\Gamma P_n = \gamma^n$.
\end{enumerate}

Proof. (1) From $C\gamma^nC^{-1} = -(\gamma^n)^T$ we have $C\gamma^n = -(\gamma^n)^TC$, so

\[
(\gamma^n)^TC\gamma^n = -\big((\gamma^n)^T\big)^2C = -\big((\gamma^n)^2\big)^TC = -\eta_{nn}C .
\]

(2) $\gamma^n$ anticommutes with $\gamma^a$ for $a \neq n$ and commutes with itself; the second identity follows by inserting the first into $S^{ab} = \frac14[\gamma^a,\gamma^b]$. (3) $\Gamma\gamma^n$ is, up to the sign of the reordering, the product of the eight gammas with $\gamma^n$ moved to the end, i.e. $\pm\eta_{nn}$ times the product of the other seven; each $\gamma^a$ is a unit vector ($(\gamma^a)^2 = \eta_{aa} = \pm1$), so the product lies in Pin(4,4). Using (2), $(\gamma^n)^{-1} = \eta_{nn}\gamma^n$, Lemma 1 and (1),

\[
\begin{aligned}
P_n^{-1}\gamma^aP_n &= (\gamma^n)^{-1}\Gamma\gamma^a\Gamma\gamma^n = -(\gamma^n)^{-1}\gamma^a\gamma^n = (R_n)^a{}_b\gamma^b ,\\
P_n^\dagger CP_n &= (\gamma^n)^T\Gamma C\Gamma\gamma^n = (\gamma^n)^TC\gamma^n = -\eta_{nn}C ,
\end{aligned}
\]

and $\Gamma P_n = \Gamma^2\gamma^n = \gamma^n$. QED. The signs of $P_n$ relative to the product of the other seven gammas (in the order of $\Gamma$) are $+1, -1, +1, +1, -1, +1, -1, -1$ for $n = x_1, \dots, x_8$.

\begingroup
\small
\begin{longtable}{@{}>{\raggedright\arraybackslash}p{0.400\linewidth}>{\raggedright\arraybackslash}p{0.400\linewidth}@{}}
\toprule
\textbf{Wolfram records} & \textbf{sympy records} \\
\midrule
\endfirsthead
\toprule
\textbf{Wolfram records} & \textbf{sympy records} \\
\midrule
\endhead
\texttt{T2\_Pn\_in\allowbreak{}\_Pin44}, \texttt{T2\_Pn\_co\allowbreak{}vers\_the\allowbreak{}\_reflect\allowbreak{}ion}, \texttt{T2\_chara\allowbreak{}cter\_of\_\allowbreak{}Pn}, \texttt{T2\_Gamma\allowbreak{}\_times\_P\allowbreak{}n\_is\_gam\allowbreak{}ma\_n}, \texttt{T2\_kerne\allowbreak{}ls\_gamma\allowbreak{}\_n\_with\_\allowbreak{}frame\_re\allowbreak{}flection} & \texttt{T2.gener\allowbreak{}al\_field\allowbreak{}.matrix\_\allowbreak{}identiti\allowbreak{}es}, \texttt{T2.gener\allowbreak{}al\_field\allowbreak{}.reflect\allowbreak{}ion\_tabl\allowbreak{}e}, \texttt{compare.\allowbreak{}theory.r\allowbreak{}eflectio\allowbreak{}n\_table} \\
\bottomrule
\end{longtable}
\endgroup

\textbf{Lemma 4 (Krein signs of the maps).} $MBM^\dagger = \sigma_MB$ with $\sigma_\Gamma = -1$; $\sigma_{\gamma^n} = +1$ for $n = x_1, x_2, x_3, x_4, x_8$ and $-1$ for $n = x_5, x_6, x_7$; and $\sigma_{P_n} = -\sigma_{\gamma^n}$, i.e. $-1$ for $P_{x_1}, P_{x_2}, P_{x_3}, P_{x_4}, P_{x_8}$ and $+1$ for $P_{x_5}, P_{x_6}, P_{x_7}$.

Proof. $\Gamma$: Lemma 2. For $\gamma^n$: it is real, so $(\gamma^n)^\dagger = (\gamma^n)^T = \eta_{nn}\gamma^n$, and transposing $C\gamma^nC^{-1} = -(\gamma^n)^T$ with $C^T = C = C^{-1}$ gives $\gamma^nC = -C(\gamma^n)^T$. Hence

\[
\gamma^nB(\gamma^n)^\dagger = -i\gamma^nC\gamma^4(\gamma^n)^T = iC(\gamma^n)^T\gamma^4(\gamma^n)^T = iC\gamma^n\gamma^4\gamma^n .
\]

For $n \neq 4$, $\gamma^n\gamma^4\gamma^n = -\eta_{nn}\gamma^4$ and the result is $\eta_{nn}B$; for $n = 4$, $\gamma^4\gamma^4\gamma^4 = -\gamma^4$ and the result is $B$. Finally $P_nBP_n^\dagger = \Gamma\,\gamma^nB(\gamma^n)^\dagger\,\Gamma = -\sigma_{\gamma^n}B$ by Lemma 2. QED.

\begingroup
\small
\begin{longtable}{@{}>{\raggedright\arraybackslash}p{0.400\linewidth}>{\raggedright\arraybackslash}p{0.400\linewidth}@{}}
\toprule
\textbf{Wolfram records} & \textbf{sympy records} \\
\midrule
\endfirsthead
\toprule
\textbf{Wolfram records} & \textbf{sympy records} \\
\midrule
\endhead
\texttt{Q\_Krein\_\allowbreak{}metric\_o\allowbreak{}f\_images} & \texttt{compare.\allowbreak{}theory.k\allowbreak{}rein\_sig\allowbreak{}ns}, \texttt{Q.image\_\allowbreak{}krein\_me\allowbreak{}tric}, \texttt{Q.T2\_ima\allowbreak{}ge\_keeps\allowbreak{}\_B} \\
\bottomrule
\end{longtable}
\endgroup

\section{Theorem T1: the chirality pairing (both fields)}

\subsection{Hypotheses}

\begin{itemize}
\item (H1) The gravitational field: any vielbein $e^a{}_\mu$ (metric $g = e^T\eta e$ of signature (4,4)) and any spin connection $\omega_{\mu ab} = -\omega_{\mu ba}$, in particular the canonical one of the author's metric. The field is a fixed (test) background: it is NOT varied, and both members of the pair are evaluated in the SAME field.
\item (H2) The statistics: $\Psi$ Grassmann-valued (dirac16complex) or commuting (dirac16complex00). The map acts at the same point, $\Psi'(x) = \Gamma\Psi(x)$, with no change of coordinates.
\item (H3) The potential: $U(S) = \frac{\lambda}{2}S^2$; more generally any function $U$ (commuting field) or any polynomial $U$ (Grassmann field), mapped to $-U$.
\item (H4) $\mathcal{L}$, $E$, $T$ and $J$ as in section 2, with any overall sign or constant factor of $T$ and $J$.
\end{itemize}

\subsection{Statement}

\textbf{Theorem T1.} Under (H1) to (H4), for both fields:

\begin{itemize}
\item (T1a) $\mathcal{L}_{m,\lambda}[\Gamma\Psi] = -\mathcal{L}_{-m,-\lambda}[\Psi]$ identically (off shell); for a general potential $\mathcal{L}_{m,U}[\Gamma\Psi] = -\mathcal{L}_{-m,-U}[\Psi]$. Separately, $S[\Gamma\Psi] = S[\Psi]$ and $K[\Gamma\Psi] = -K[\Psi]$.
\item (T1b) $E_{m,\lambda}[\Gamma\Psi] = -\Gamma E_{-m,-\lambda}[\Psi]$: $\Gamma\Psi$ solves the field equations with $(m,\lambda)$ if and only if $\Psi$ solves those with $(-m,-\lambda)$; equivalently, $\Psi$ solves the $(m,\lambda)$ equations if and only if $\Gamma\Psi$ solves the $(-m,-\lambda)$ equations. The Euler-Lagrange expressions are mapped in the same way.
\item (T1c) $T^{(-m,-\lambda)}_{\mu\nu}[\Gamma\Psi] = -T^{(m,\lambda)}_{\mu\nu}[\Psi]$ and $J^\mu[\Gamma\Psi] = -J^\mu[\Psi]$ pointwise, off and on shell.
\item (T1d) The pair ($\Psi$ with $(m,\lambda)$; $\Gamma\Psi$ with $(-m,-\lambda)$) has the total energy-momentum density $T + T' = 0$, the total current $J + J' = 0$ and the total charge $Q + Q' = 0$, as classical bilinears (c-numbers for dirac16complex00, elements of the Grassmann algebra for dirac16complex).
\end{itemize}

\subsection{Proof}

Step 1 (bilinears). By Lemma 1, $D_\mu(\Gamma\Psi) = \Gamma D_\mu\Psi$ and $D_\mu(\overline{\Gamma\Psi}) = (D_\mu\bar\Psi)\Gamma$; by Lemma 2, $S[\Gamma\Psi] = S[\Psi]$ and

\[
\overline{\Gamma\Psi}\,\gamma^\mu D_\mu(\Gamma\Psi) = \bar\Psi\,\Gamma\gamma^\mu\Gamma\,D_\mu\Psi = -\bar\Psi\gamma^\mu D_\mu\Psi ,
\qquad
\big(D_\mu\overline{\Gamma\Psi}\big)\gamma^\mu\,\Gamma\Psi = -(D_\mu\bar\Psi)\gamma^\mu\Psi ,
\]

hence $K[\Gamma\Psi] = -K[\Psi]$. Written out,

\[
K = \tfrac12\,e^\mu{}_a\big(\bar\Psi\gamma^a\partial_\mu\Psi - \partial_\mu\bar\Psi\gamma^a\Psi\big) + \tfrac14\,e^\mu{}_a\,\omega_{\mu bc}\,\bar\Psi\{\gamma^a,S^{bc}\}\Psi ,
\]

a sum of coefficient functions of the gravitational field times bilinears with an odd number of gammas; each such bilinear changes sign by Lemma 2, whatever the coefficients are.

Step 2 (the Lagrangian). $U(S)$ is unchanged because $S$ is, so

\[
\mathcal{L}_{m,U}[\Gamma\Psi] = \sqrt{|g|}\,\big[-K - mS - U(S)\big] = -\sqrt{|g|}\,\big[K - (-m)S - (-U)(S)\big] = -\mathcal{L}_{-m,-U}[\Psi] ,
\]

which for $U = \frac{\lambda}{2}S^2$ is (T1a).

Step 3 (the field equations). $\gamma^\mu D_\mu(\Gamma\Psi) = \gamma^\mu\Gamma D_\mu\Psi = -\Gamma\gamma^\mu D_\mu\Psi$, so

\[
E_{m,\lambda}[\Gamma\Psi] = -\Gamma\gamma^\mu D_\mu\Psi - (m + \lambda S)\Gamma\Psi = -\Gamma\big[\gamma^\mu D_\mu\Psi - (-m - \lambda S)\Psi\big] = -\Gamma E_{-m,-\lambda}[\Psi] ,
\]

and $\Gamma$ is invertible, which gives (T1b). For the Euler-Lagrange expressions themselves, in any field and whatever their explicit form: (T1a) is an identity between functionals of $\Psi$, and $\Psi \mapsto \Gamma\Psi$ is a constant invertible linear substitution, so by the chain rule the Euler-Lagrange expressions of $\mathcal{L}_{m,\lambda}$ at $\Gamma\Psi$, multiplied by $\Gamma$, equal minus those of $\mathcal{L}_{-m,-\lambda}$ at $\Psi$; they vanish together.

Step 4 (energy-momentum tensor and current). $T^\nu{}_\mu$ is the sum of $\delta^\nu_\mu\mathcal{L}/\sqrt{|g|}$ and kinetic bilinears with one gamma. By Step 2 with $m \to -m$, $\mathcal{L}_{-m,-\lambda}[\Gamma\Psi] = -\mathcal{L}_{m,\lambda}[\Psi]$, and by Step 1 each kinetic bilinear changes sign; hence $T^{(-m,-\lambda)}[\Gamma\Psi] = -T^{(m,\lambda)}[\Psi]$. The same holds for the tensor of the vielbein variation, whose additional term $\nabla_\lambda(\bar\Psi\{\gamma^\lambda,\Sigma^{\nu\rho}\}\Psi)$ contains three gammas. $J^\mu$ is a multiple of $\bar\Psi\gamma^\mu\Psi$ and changes sign by Lemma 2, and so does $Q = \int\sqrt{|g|}\,J^{x_4}\,d^7x$. This is (T1c).

Step 5 (the pair). Adding (T1c) to the quantities of $\Psi$ gives (T1d).

Step 6 (statistics). $\Gamma = \mathrm{diag}(-I_8, I_8)$ multiplies each component by $+1$ or $-1$; no Grassmann factors are reordered, so Steps 1 to 5 hold verbatim for anticommuting components, as identities between Grassmann-algebra elements.

The proof uses only Lemmas 1 and 2 and the linearity of $\mathcal{L}$, $T$ and $J$ in the coefficient functions $\sqrt{|g|}$, $e^\mu{}_a$ and $\omega_{\mu ab}$; hence T1 holds in every gravitational field (H1). QED.

\subsection{T1 in the block form of the author's field}

In the author's field the proof can be read off the block equations of section 2.5. $\Gamma(\psi_-, \psi_+) = (-\psi_-, \psi_+)$, and $C$ is block diagonal (a product of four gammas; the algebra record gives $C = \mathrm{diag}(-\sigma, \sigma)$ with $\sigma$ the notebook's $8 \times 8$ involution), so $S = -\psi_-^\dagger\sigma\psi_- + \psi_+^\dagger\sigma\psi_+$ is unchanged and $V = m + \lambda S$ becomes $-m - \lambda S = -V$ for the parameters $(-m,-\lambda)$. The two block equations become

\[
\sum_a f_a^{-1}\bar\tau^a\partial_a\psi_+ + 3H\psi_+ = (-V)(-\psi_-),\qquad
-\Big(\sum_a f_a^{-1}\tau^a\partial_a\psi_- + 3H\psi_-\Big) = -V\psi_+ ,
\]

the same equations. None of the gravitational terms ($e^{\mp a_4}\sin^{-1/6}z$, $\tan z$, $3H$) changes.

\begingroup
\small
\begin{longtable}{@{}>{\raggedright\arraybackslash}p{0.400\linewidth}>{\raggedright\arraybackslash}p{0.400\linewidth}@{}}
\toprule
\textbf{Wolfram records (author's field)} & \textbf{sympy records (author's field)} \\
\midrule
\endfirsthead
\toprule
\textbf{Wolfram records (author's field)} & \textbf{sympy records (author's field)} \\
\midrule
\endhead
\texttt{T1\_Lagra\allowbreak{}ngian\_pr\allowbreak{}imordial\allowbreak{}\_commuti\allowbreak{}ng} (408 monomials), \texttt{T1\_Lagra\allowbreak{}ngian\_pr\allowbreak{}imordial\allowbreak{}\_grassma\allowbreak{}nn} (392 monomials), \texttt{T1\_field\allowbreak{}\_equatio\allowbreak{}n\_covari\allowbreak{}ance\_pri\allowbreak{}mordial\_\allowbreak{}commutin\allowbreak{}g}, \texttt{T1\_field\allowbreak{}\_equatio\allowbreak{}n\_covari\allowbreak{}ance\_pri\allowbreak{}mordial\_\allowbreak{}grassman\allowbreak{}n}, \texttt{T1\_energ\allowbreak{}y\_moment\allowbreak{}um\_primo\allowbreak{}rdial\_co\allowbreak{}mmuting}, \texttt{T1\_energ\allowbreak{}y\_moment\allowbreak{}um\_primo\allowbreak{}rdial\_gr\allowbreak{}assmann}, \texttt{T1\_curre\allowbreak{}nt\_primo\allowbreak{}rdial\_co\allowbreak{}mmuting}, \texttt{T1\_curre\allowbreak{}nt\_primo\allowbreak{}rdial\_gr\allowbreak{}assmann} & \texttt{C\_equals\allowbreak{}\_noteboo\allowbreak{}k\_sigma1\allowbreak{}6}, \texttt{T1.metri\allowbreak{}c.commut\allowbreak{}ing.lagr\allowbreak{}angian}, \texttt{T1.metri\allowbreak{}c.grassm\allowbreak{}ann.lagr\allowbreak{}angian}, \texttt{T1.metri\allowbreak{}c.commut\allowbreak{}ing.eule\allowbreak{}r\_lagran\allowbreak{}ge\_map}, \texttt{T1.metri\allowbreak{}c.grassm\allowbreak{}ann.eule\allowbreak{}r\_lagran\allowbreak{}ge\_map}, \texttt{T1.metri\allowbreak{}c.commut\allowbreak{}ing.pair\allowbreak{}\_total\_e\allowbreak{}mt\_zero}, \texttt{T1.metri\allowbreak{}c.grassm\allowbreak{}ann.pair\allowbreak{}\_total\_e\allowbreak{}mt\_zero} \\
\bottomrule
\end{longtable}
\endgroup

\subsection{Remarks on T1}

\begin{itemize}
\item Negative controls: $\mathcal{L}_{m,\lambda}[\Gamma\Psi] + \mathcal{L}_{m,\lambda}[\Psi] \neq 0$ and $\mathcal{L}_{m,\lambda}[\Gamma\Psi] + \mathcal{L}_{-m,\lambda}[\Psi] \neq 0$: the mass AND the coupling must change sign. For $\lambda \neq 0$, T1 is therefore not a pure $+m$ / $-m$ pairing; for $\lambda = 0$ (free fields) it pairs $+m$ with $-m$ exactly.
\item The partner's energy density, pressures and equations of state: $\rho' = -\rho$ and $p_\mu' = -p_\mu$ in every direction; for homogeneous on-shell states $\rho = mS + \frac{\lambda}{2}S^2$ and $p = \frac{\lambda}{2}S^2$ become $\rho' = -mS - \frac{\lambda}{2}S^2$ and $p' = -\frac{\lambda}{2}S^2$ (with $S' = S$). Every ratio of components, in particular $w_3 = p_3/\rho$, $w_t = p_t/\rho$ and $w_8 = p_8/\rho$, is the same for both members.
\item T1 is not a symmetry of one theory: it changes the parameters and the sign of the action. Because $-\mathcal{L}$ and $\mathcal{L}$ have the same Euler-Lagrange equations, $\Gamma\Psi$ is a solution of the $(-m,-\lambda)$ theory; nothing more is asserted.
\item Representation content: $\Gamma$ is $-1$ on one irreducible Spin(4,4) half and $+1$ on the other (section 2.2); the partner differs from $\Psi$ by the relative sign of its two inequivalent Spin(4,4) components, and $\Gamma$ is itself an element of Pin(4,4) (a product of eight unit vectors).
\end{itemize}

\begingroup
\small
\begin{longtable}{@{}>{\raggedright\arraybackslash}p{0.400\linewidth}>{\raggedright\arraybackslash}p{0.400\linewidth}@{}}
\toprule
\textbf{Wolfram records (negative controls)} & \textbf{sympy records (negative controls)} \\
\midrule
\endfirsthead
\toprule
\textbf{Wolfram records (negative controls)} & \textbf{sympy records (negative controls)} \\
\midrule
\endhead
\texttt{T1\_Lagra\allowbreak{}ngian\_ne\allowbreak{}gative\_c\allowbreak{}ontrols\_\allowbreak{}primordi\allowbreak{}al\_commu\allowbreak{}ting}, \texttt{T1\_Lagra\allowbreak{}ngian\_ne\allowbreak{}gative\_c\allowbreak{}ontrols\_\allowbreak{}primordi\allowbreak{}al\_grass\allowbreak{}mann}, \texttt{T1\_Lagra\allowbreak{}ngian\_ne\allowbreak{}gative\_c\allowbreak{}ontrols\_\allowbreak{}diagonal\allowbreak{}8\_commut\allowbreak{}ing}, \texttt{T1\_Lagra\allowbreak{}ngian\_ne\allowbreak{}gative\_c\allowbreak{}ontrols\_\allowbreak{}diagonal\allowbreak{}8\_grassm\allowbreak{}ann}, \texttt{T1\_Lagra\allowbreak{}ngian\_ne\allowbreak{}gative\_c\allowbreak{}ontrols\_\allowbreak{}pointwis\allowbreak{}e\_commut\allowbreak{}ing}, \texttt{T1\_Lagra\allowbreak{}ngian\_ne\allowbreak{}gative\_c\allowbreak{}ontrols\_\allowbreak{}pointwis\allowbreak{}e\_grassm\allowbreak{}ann} & \texttt{T1.metri\allowbreak{}c.commut\allowbreak{}ing.nega\allowbreak{}tive\_con\allowbreak{}trols}, \texttt{T1.metri\allowbreak{}c.grassm\allowbreak{}ann.nega\allowbreak{}tive\_con\allowbreak{}trols} \\
\bottomrule
\end{longtable}
\endgroup

\section{Theorem T2: the mirror pairing (both fields)}

\subsection{Hypotheses}

\begin{itemize}
\item (H5) $n$ is a space-like frame direction ($x_1$, $x_2$, $x_3$ or $x_8$), and $P_n = \Gamma\gamma^n$ is the Pin(4,4) lift of the reflection $R_n$, of character $-1$ (Lemma 3).
\item (H6) General-field version: the frame is reflected, $e' = R_ne$ (the same metric, since $R_n\eta R_n = \eta$, and the same $\sqrt{|g|}$), with its canonical connection, and $\Psi'(x) = \Gamma P_n\Psi(x) = \gamma^n\Psi(x)$ at the same point.
\item (H7) Author's-field version: the map $\phi: x_8 \mapsto \pi/(6H) - x_8$ ($z \mapsto \pi - z$) takes the patch $z \in (0,\pi/2)$ onto the mirror patch $z \in (\pi/2,\pi)$; each patch carries its own positive vielbein ($e^8 = |\cot z|\,dx_8$, i.e. $-\cot z\,dx_8$ on the mirror patch); and $\Psi'(\phi(x)) = \gamma^8\Psi(x)$. The brane $z = \pi/2$ is a degenerate surface of the metric ($g_{88} = \cot^2 z = 0$ and $\sqrt{|g|} = \cos z = 0$ there); extending the field to the mirror patch is the ASSUMED Z2 construction of \texttt{Revision\allowbreak{}/SPEC.md} section 7.
\item (H8) Both statistics; $U(S) = \frac{\lambda}{2}S^2$ or any even function of $S$.
\end{itemize}

\subsection{Statement}

\textbf{Theorem T2.} Under (H5) to (H8), for both fields:

\begin{itemize}
\item (T2a) $\mathcal{L}_{m,\lambda}[\gamma^n\Psi;\,R_ne] = +\mathcal{L}_{-m,\lambda}[\Psi;\,e]$ in every gravitational field; in the author's field $\mathcal{L}_{m,\lambda}[\Psi'](\phi(x)) = \mathcal{L}_{-m,\lambda}[\Psi](x)$.
\item (T2b) $S' = -S$, and $E_{m,\lambda}[\gamma^n\Psi;\,R_ne] = -\gamma^nE_{-m,\lambda}[\Psi;\,e]$: $\Psi$ solves the $(-m,\lambda)$ equations if and only if its image solves the $(m,\lambda)$ equations (in the author's field: on the patch and on the mirror patch respectively).
\item (T2c) In the general-field version $T'_{\mu\nu} = T^{(-m,\lambda)}_{\mu\nu}[\Psi]$ and $J'^\mu = J^\mu[\Psi]$. In the author's field $T_{\mu\nu}[\Psi'](\phi(x)) = (R_8\,T^{(-m,\lambda)}[\Psi](x)\,R_8)_{\mu\nu}$ and $J'^\mu = (R_8)^\mu{}_\nu J^\nu$: the energy density and every component without an $x_8$ index are EQUAL, not opposite. T2 pairs $+m$ with $-m$ at equal energy-momentum.
\item (T2d) For a time-like $n$ (character $+1$) the same construction gives $\mathcal{L} \to -\mathcal{L}_{-m,-\lambda}$, a T1-type map; the character $-1$ is what makes T2 a pairing $(m,\lambda) \to (-m,\lambda)$ with $\mathcal{L} \to +\mathcal{L}$.
\end{itemize}

\subsection{Proof in a general field (frame reflection)}

Step 1 (geometry). $g' = e'^T\eta e' = e^TR_n\eta R_ne = g$. The canonical connection of $e'$ is $R_n\omega R_n$: $R_n$ is a constant Lorentz transformation, and the canonical connection $\omega_\mu{}^a{}_b = e^a{}_\nu(\partial_\mu e^\nu{}_b + \Gamma^\nu{}_{\mu\lambda}e^\lambda{}_b)$ is built from the vielbein, its inverse and the Christoffel symbols $\Gamma^\nu{}_{\mu\lambda}$ of the unchanged metric (the symbol with indices is the Christoffel symbol, without indices the chirality matrix). By Lemma 3, $\Omega'_\mu = \frac12(R_n\omega_\mu R_n)_{ab}S^{ab} = \gamma^n\Omega_\mu(\gamma^n)^{-1}$, so $D'_\mu(\gamma^n\Psi) = \gamma^nD_\mu\Psi$.

Step 2 (adjoint). By Lemma 3, $(\gamma^n)^TC = -\eta_{nn}C(\gamma^n)^{-1}$, hence $\bar\Psi' = \Psi^\dagger(\gamma^n)^TC = -\eta_{nn}\bar\Psi(\gamma^n)^{-1}$ and $D'_\mu\bar\Psi' = -\eta_{nn}(D_\mu\bar\Psi)(\gamma^n)^{-1}$.

Step 3 (gammas). The reflected frame gives $\gamma'^\mu = e^\mu{}_b(R_n)^b{}_a\gamma^a$, and Lemma 3 gives $\gamma'^\mu\gamma^n = -\gamma^n\gamma^\mu$. Therefore

\[
\bar\Psi'\gamma'^\mu D'_\mu\Psi' = -\eta_{nn}\bar\Psi(\gamma^n)^{-1}\gamma'^\mu\gamma^nD_\mu\Psi = \eta_{nn}\bar\Psi\gamma^\mu D_\mu\Psi ,
\qquad
\bar\Psi'\Psi' = -\eta_{nn}\bar\Psi\Psi ,
\]

and in the same way $(D'_\mu\bar\Psi')\gamma'^\mu\Psi' = \eta_{nn}(D_\mu\bar\Psi)\gamma^\mu\Psi$, so $K' = \eta_{nn}K$ and $S' = -\eta_{nn}S$.

Step 4 (Lagrangian). For space-like $n$ ($\eta_{nn} = +1$): $K' = K$ and $S' = -S$, so

\[
\mathcal{L}_{m,\lambda}[\Psi';e'] = \sqrt{|g|}\,\big[K + mS - \tfrac{\lambda}{2}S^2\big] = \mathcal{L}_{-m,\lambda}[\Psi;e] ,
\]

which is (T2a). For time-like $n$: $K' = -K$ and $S' = S$, so $\mathcal{L}_{m,\lambda}[\Psi';e'] = -\mathcal{L}_{-m,-\lambda}[\Psi;e]$, which is (T2d).

Step 5 (field equations). For space-like $n$,

\[
E_{m,\lambda}[\Psi';e'] = \gamma'^\mu D'_\mu(\gamma^n\Psi) - (m + \lambda S')\gamma^n\Psi = -\gamma^n\gamma^\mu D_\mu\Psi - (m - \lambda S)\gamma^n\Psi = -\gamma^nE_{-m,\lambda}[\Psi;e] ,
\]

which is (T2b).

Step 6 (energy-momentum tensor and current). The kinetic bilinears $\bar\Psi\gamma_\mu D_\nu\Psi$ with free indices transform like the contracted one in Step 3 (factor $\eta_{nn} = +1$), and $\mathcal{L}/\sqrt{|g|}$ transforms by Step 4, so $T'_{\mu\nu} = T^{(-m,\lambda)}_{\mu\nu}[\Psi]$; $\bar\Psi'\gamma'^\mu\Psi' = \eta_{nn}\bar\Psi\gamma^\mu\Psi$ gives $J' = J$. This is (T2c) in the general-field version.

Step 7 (statistics). $\gamma^n$ is a constant signed permutation matrix; the substitution never reorders Grassmann factors, so Steps 1 to 6 hold for both statistics. QED.

\subsection{Proof in the author's field (the Z2 mirror)}

$\phi$ is an isometry of the author's metric: its components depend on $x_8$ only through $\sin z$ and $\cot^2 z$, both invariant under $z \mapsto \pi - z$. The positive vielbein of the mirror patch, pulled back by $\phi$, is $R_8e$: its hidden leg is $-\cot(\pi - z)\,d(\pi/(6H) - x_8) = -\cot z\,dx_8 = -e^8$, and the other legs are unchanged. $\sqrt{|g|}\,d^8x$ is invariant ($|\cos z|$ is invariant and the Jacobian has modulus 1). Pulling the mirror-patch Lagrangian back by $\phi$ therefore gives the frame-reflected Lagrangian of section 5.3 with $n = x_8$, which proves (T2a) and (T2b) in the author's field; the coordinate components of $T$ and $J$ pick up the Jacobian $R_8$ of $\phi$, which gives (T2c). QED.

In block form $\gamma^8$ has the blocks $\bar\tau^8 = \tau^8 = I_8$, so $\gamma^8(\psi_-, \psi_+) = (\psi_+, \psi_-)$: the mirror image exchanges the two inequivalent Spin(4,4) halves, and $S = -\psi_-^\dagger\sigma\psi_- + \psi_+^\dagger\sigma\psi_+$ changes sign.

\begingroup
\small
\begin{longtable}{@{}>{\raggedright\arraybackslash}p{0.400\linewidth}>{\raggedright\arraybackslash}p{0.400\linewidth}@{}}
\toprule
\textbf{Wolfram records (author's field)} & \textbf{sympy records (author's field)} \\
\midrule
\endfirsthead
\toprule
\textbf{Wolfram records (author's field)} & \textbf{sympy records (author's field)} \\
\midrule
\endhead
\texttt{T2\_mirro\allowbreak{}r\_is\_iso\allowbreak{}metry}, \texttt{T2\_mirro\allowbreak{}r\_Lagran\allowbreak{}gian\_com\allowbreak{}muting}, \texttt{T2\_mirro\allowbreak{}r\_Lagran\allowbreak{}gian\_gra\allowbreak{}ssmann}, \texttt{T2\_mirro\allowbreak{}r\_energy\allowbreak{}\_momentu\allowbreak{}m\_and\_cu\allowbreak{}rrent\_co\allowbreak{}mmuting}, \texttt{T2\_mirro\allowbreak{}r\_energy\allowbreak{}\_momentu\allowbreak{}m\_and\_cu\allowbreak{}rrent\_gr\allowbreak{}assmann}, \texttt{T2\_mirro\allowbreak{}r\_Euler\_\allowbreak{}Lagrange\allowbreak{}\_commuti\allowbreak{}ng}, \texttt{T2\_mirro\allowbreak{}r\_Euler\_\allowbreak{}Lagrange\allowbreak{}\_grassma\allowbreak{}nn} & \texttt{geometry\allowbreak{}.mirror\_\allowbreak{}isometry}, \texttt{geometry\allowbreak{}.brane\_d\allowbreak{}egenerat\allowbreak{}e}, \texttt{T2.metri\allowbreak{}c.commut\allowbreak{}ing.eule\allowbreak{}r\_lagran\allowbreak{}ge\_map}, \texttt{T2.metri\allowbreak{}c.grassm\allowbreak{}ann.eule\allowbreak{}r\_lagran\allowbreak{}ge\_map}, \texttt{T2.metri\allowbreak{}c.commut\allowbreak{}ing.emt}, \texttt{T2.metri\allowbreak{}c.grassm\allowbreak{}ann.emt}, \texttt{T2.metri\allowbreak{}c.commut\allowbreak{}ing.curr\allowbreak{}ent}, \texttt{T2.metri\allowbreak{}c.grassm\allowbreak{}ann.curr\allowbreak{}ent}, \texttt{T2.metri\allowbreak{}c.commut\allowbreak{}ing.S\_od\allowbreak{}d}, \texttt{T2.metri\allowbreak{}c.grassm\allowbreak{}ann.S\_od\allowbreak{}d} \\
\bottomrule
\end{longtable}
\endgroup

\subsection{The eight reflections}

The sympy checker also shows that neither the plain mirror ($P = 1$) nor $\Gamma$ alone gives a relation $\mathcal{L}_{\mathrm{mirror}} = \pm\mathcal{L}_{\mathrm{patch}}$ with $(m,\lambda) \to (\pm m, \pm\lambda)$, that $P_8 = \Gamma\gamma^8$ alone gives $\mathcal{L} \to -\mathcal{L}_{m,-\lambda}$, and that its composition with $\Gamma$ gives T2. The frame reflections of all eight directions, from the data of the pairing theory record (the Wolfram pointwise frame-reflection checks of every direction and the sympy frame-reflection checks of the general diagonal field, both statistics; the record names are in the T2 tables of section 9.2):

\begingroup
\small
\begin{longtable}{@{}>{\raggedright\arraybackslash}p{0.133\linewidth}>{\raggedright\arraybackslash}p{0.133\linewidth}>{\raggedright\arraybackslash}p{0.133\linewidth}>{\raggedright\arraybackslash}p{0.133\linewidth}>{\raggedright\arraybackslash}p{0.133\linewidth}>{\raggedright\arraybackslash}p{0.133\linewidth}@{}}
\toprule
\textbf{direction} & \textbf{$\eta_{nn}$} & \textbf{character of $P_n$} & \textbf{$S$ under $\gamma^n$} & \textbf{kinetic term} & \textbf{map} \\
\midrule
\endfirsthead
\toprule
\textbf{direction} & \textbf{$\eta_{nn}$} & \textbf{character of $P_n$} & \textbf{$S$ under $\gamma^n$} & \textbf{kinetic term} & \textbf{map} \\
\midrule
\endhead
$x_1$ & $+1$ & $-1$ & $-S$ & $+K$ & $(-m,\lambda)$, $+\mathcal{L}$ (T2) \\
$x_2$ & $+1$ & $-1$ & $-S$ & $+K$ & $(-m,\lambda)$, $+\mathcal{L}$ (T2) \\
$x_3$ & $+1$ & $-1$ & $-S$ & $+K$ & $(-m,\lambda)$, $+\mathcal{L}$ (T2) \\
$x_4$ & $-1$ & $+1$ & $+S$ & $-K$ & $(-m,-\lambda)$, $-\mathcal{L}$ (T1-type) \\
$x_5$ & $-1$ & $+1$ & $+S$ & $-K$ & $(-m,-\lambda)$, $-\mathcal{L}$ (T1-type) \\
$x_6$ & $-1$ & $+1$ & $+S$ & $-K$ & $(-m,-\lambda)$, $-\mathcal{L}$ (T1-type) \\
$x_7$ & $-1$ & $+1$ & $+S$ & $-K$ & $(-m,-\lambda)$, $-\mathcal{L}$ (T1-type) \\
$x_8$ & $+1$ & $-1$ & $-S$ & $+K$ & $(-m,\lambda)$, $+\mathcal{L}$ (T2) \\
\bottomrule
\end{longtable}
\endgroup

In the author's field only the reflection of $x_8$ is realised by an isometry that maps the patch onto another patch (the mirror across the Z2 brane); the reflections of $x_1, x_2, x_3$ are frame statements.

\section{The quantum-level reading (dirac16complex)}

\subsection{Hypotheses}

\begin{itemize}
\item (HQ1) dirac16complex is canonically quantised with $x_4$ as the evolution time. The $x_4$-derivative kernel of $\mathcal{L}_{m,\lambda}$ is $\frac{i}{2}\sqrt{|g|}\,N$ with $N = B$ (since $C\gamma^4 = iB$), and the canonical anticommutator on a slice $x_4 = \mathrm{const}$ is $\{\Psi_A(x), \Psi^\dagger_B(y)\} = (N^{-1})_{AB}\,\delta^7(x - y)/\sqrt{|g|} = B_{AB}\,\delta^7(x - y)/\sqrt{|g|}$.
\item (HQ2) dirac16complex00 is a classical field and has no quantum reading.
\end{itemize}

\subsection{Statements}

\begin{itemize}
\item (Q1) The chirality image $\chi = \Gamma\Psi$ carries the Krein metric $-B$: $\{\chi_A, \chi^\dagger_B\} = (\Gamma B\Gamma)_{AB}\,\delta/\sqrt{|g|} = -B_{AB}\,\delta/\sqrt{|g|}$, which is also the anticommutator demanded by its own Lagrangian $\mathcal{L}_{m,\lambda}[\Gamma\chi] = -\mathcal{L}_{-m,-\lambda}[\chi]$.
\item (Q2) The image's own generators (energy, momentum and charge from its own Lagrangian) coincide with those of $\Psi$: the pair $\Psi$, $\Gamma\Psi$ is ONE quantum system, and the T1 identity $T^{(-m,-\lambda)}[\Gamma\Psi] = -T^{(m,\lambda)}[\Psi]$ is an identity between operators of that one system.
\item (Q3) An independently quantised $(-m,-\lambda)$ universe has the anticommutator $+B$ (from $\mathcal{L}_{-m,-\lambda}$), its own vacuum and its own generators; it cannot be identified with $\Gamma\Psi$ ($B \neq -B$); on the product state space the generators add, $P_{\mathrm{total}} = P_1\otimes 1 + 1\otimes P_2$, and no cancellation $P_1 + P_2 = 0$ follows.
\item (Q4) The one-particle Hamiltonians satisfy $\Gamma h_m\Gamma = h_{-m}$ (flat space, and a general field at a point) and $\gamma^8h_m(k)\gamma^8 = h_{-m}(R_8k)$: the $+m$ and $-m$ universes have IDENTICAL spectra, not opposite ones. In flat 4+4 space $h_m(k) = -im\gamma^4 - \gamma^4\sum_{a \neq 4}k_a\gamma^a$ and $h_m(k)^2 = (m^2 + k_1^2 + k_2^2 + k_3^2 + k_8^2 - k_5^2 - k_6^2 - k_7^2)\,I_{16}$; at every exactly computed sample (table in section 6.3) each energy eigenspace has dimension 8 and Krein inertia (4,4), the same for $+m$ and $-m$.
\item (Q5) The T2 image $\gamma^8\Psi$ keeps the anticommutator $+B$: the mirror $(-m,\lambda)$ universe is an ordinary, independently quantisable copy with equal energies.
\end{itemize}

\subsection{Proofs}

(Q1) By Lemma 2, $\{\Gamma\Psi, (\Gamma\Psi)^\dagger\} = \Gamma\{\Psi, \Psi^\dagger\}\Gamma^\dagger = \Gamma B\Gamma\,\delta/\sqrt{|g|} = -B\,\delta/\sqrt{|g|}$. For the image's own Lagrangian, the $x_4$-velocity term $\frac{i}{2}\sqrt{|g|}\,\chi^\dagger\Gamma B\Gamma\,\partial_4\chi$ has the kernel $N = -B$, hence the anticommutator $N^{-1} = -B$: the two agree.

(Q2) At $\chi = \Gamma\Psi$ the image Lagrangian $\mathcal{L}_{m,\lambda}[\Gamma\chi]$ is $\mathcal{L}_{m,\lambda}[\Psi]$, the same function of the same operators, so its Legendre energy density and its Noether densities coincide with those of $\Psi$; they are minus the generators of the $(-m,-\lambda)$ Lagrangian evaluated at $\chi$. With the metric $-B$ and the generator density $-h_m$ the image evolves by $(-B)(-h_m) = Bh_m$, exactly the $(m,\lambda)$ dynamics.

(Q3) The Lagrangian $\mathcal{L}_{-m,-\lambda}[\chi_2]$ of an independent universe has the kernel $N = +B$, hence $\{\chi_2, \chi_2^\dagger\} = +B\,\delta/\sqrt{|g|}$, while $\{\Gamma\Psi_1, (\Gamma\Psi_1)^\dagger\} = -B\,\delta/\sqrt{|g|}$; and independence requires $\{\Psi_1, \chi_2^\dagger\} = 0$, whereas $\{\Gamma\Psi_1, \Psi_1^\dagger\} = \Gamma B/\sqrt{|g|}$ has rank 16. So no identification $\chi_2 = \Gamma\Psi_1$ exists. On the product space the generators of the two systems add; the T1 identity relates operators of universe 1 only and gives no relation $P_2 = -P_1$. At zero momentum the block one-particle generator $\mathrm{diag}(Bh_m, Bh_{-m})$ has the eigenvalues $-m$ and $m$, each 16 times, all nonzero for $m \neq 0$.

(Q4) $\Gamma\gamma^4\Gamma = -\gamma^4$ and $\Gamma\gamma^4\gamma^a\Gamma = \gamma^4\gamma^a$ give $\Gamma h_m\Gamma = h_{-m}$. With $(\gamma^8)^2 = 1$,

\[
\gamma^8\gamma^4\gamma^8 = -\gamma^4,\qquad
\gamma^8\gamma^4\gamma^a\gamma^8 = \gamma^4\gamma^a \ (a \neq 4, 8),\qquad
\gamma^8\gamma^4\gamma^8\gamma^8 = -\gamma^4\gamma^8 ,
\]

which gives $\gamma^8h_m(k)\gamma^8 = h_{-m}(R_8k)$. Similar matrices have equal spectra. The square and the eigenspaces are computed exactly; the samples of \texttt{Revision\allowbreak{}/pairing\allowbreak{}/pairing\allowbreak{}-theory.\allowbreak{}json} are:

\begingroup
\small
\begin{longtable}{@{}>{\raggedright\arraybackslash}p{0.160\linewidth}>{\raggedright\arraybackslash}p{0.160\linewidth}>{\raggedright\arraybackslash}p{0.160\linewidth}>{\raggedright\arraybackslash}p{0.160\linewidth}>{\raggedright\arraybackslash}p{0.160\linewidth}@{}}
\toprule
\textbf{$m$} & \textbf{momentum $(k_1, \dots, k_8)$} & \textbf{$w$} & \textbf{dimensions of the $+w$ and $-w$ eigenspaces} & \textbf{Krein inertia on $+w$ and on $-w$} \\
\midrule
\endfirsthead
\toprule
\textbf{$m$} & \textbf{momentum $(k_1, \dots, k_8)$} & \textbf{$w$} & \textbf{dimensions of the $+w$ and $-w$ eigenspaces} & \textbf{Krein inertia on $+w$ and on $-w$} \\
\midrule
\endhead
$2$ & (1, 2, 0, 0, 0, 0, 0, 4) & $5$ & 8 and 8 & (4,4) and (4,4) \\
$-2$ & (1, 2, 0, 0, 0, 0, 0, 4) & $5$ & 8 and 8 & (4,4) and (4,4) \\
$3$ & (0, 0, 0, 0, 0, 0, 0, 0) & $3$ & 8 and 8 & (4,4) and (4,4) \\
$-3$ & (0, 0, 0, 0, 0, 0, 0, 0) & $3$ & 8 and 8 & (4,4) and (4,4) \\
$1$ & (1, 0, 0, 0, 0, 0, 0, 0) & $\sqrt{2}$ & 8 and 8 & (4,4) and (4,4) \\
$-1$ & (1, 0, 0, 0, 0, 0, 0, 0) & $\sqrt{2}$ & 8 and 8 & (4,4) and (4,4) \\
\bottomrule
\end{longtable}
\endgroup

The entry $k_4$ is not used ($h_m$ contains no $\gamma^4\gamma^4$ term). In a general field at a point the same similarity $\Gamma h_m\Gamma = h_{-m}$ holds for symbolic momenta; in the pointwise test field $(\gamma^{x_4})^2 = g^{x_4x_4}I_{16} = -\frac{2269482}{1990921}I_{16}$.

(Q5) Lemma 4 gives $\gamma^8B(\gamma^8)^\dagger = +B$. QED.

\begingroup
\small
\begin{longtable}{@{}>{\raggedright\arraybackslash}p{0.267\linewidth}>{\raggedright\arraybackslash}p{0.267\linewidth}>{\raggedright\arraybackslash}p{0.267\linewidth}@{}}
\toprule
\textbf{statement} & \textbf{Wolfram records} & \textbf{sympy records} \\
\midrule
\endfirsthead
\toprule
\textbf{statement} & \textbf{Wolfram records} & \textbf{sympy records} \\
\midrule
\endhead
(HQ1) & \texttt{first\_or\allowbreak{}der\_form\allowbreak{}\_and\_ant\allowbreak{}icommuta\allowbreak{}tor} & \texttt{canonica\allowbreak{}l\_antico\allowbreak{}mmutator\allowbreak{}\_B}, \texttt{Q.canoni\allowbreak{}cal\_anti\allowbreak{}commutat\allowbreak{}or} \\
(Q1) & \texttt{Q\_Krein\_\allowbreak{}metric\_o\allowbreak{}f\_images}, \texttt{Q\_symple\allowbreak{}ctic\_ker\allowbreak{}nel\_comm\allowbreak{}uting}, \texttt{Q\_symple\allowbreak{}ctic\_ker\allowbreak{}nel\_gras\allowbreak{}smann} & \texttt{Q.image\_\allowbreak{}krein\_me\allowbreak{}tric}, \texttt{Q.image\_\allowbreak{}own\_quan\allowbreak{}tisation} \\
(Q2) & \texttt{Q\_genera\allowbreak{}tors\_of\_\allowbreak{}the\_imag\allowbreak{}e\_commut\allowbreak{}ing}, \texttt{Q\_genera\allowbreak{}tors\_of\_\allowbreak{}the\_imag\allowbreak{}e\_grassm\allowbreak{}ann} & \texttt{Q.image\_\allowbreak{}generato\allowbreak{}rs\_same\_\allowbreak{}dynamics} \\
(Q3) & \texttt{Q\_no\_ide\allowbreak{}ntificat\allowbreak{}ion\_of\_i\allowbreak{}ndepende\allowbreak{}nt\_unive\allowbreak{}rses} & \texttt{Q.no\_can\allowbreak{}cellatio\allowbreak{}n\_indepe\allowbreak{}ndent\_un\allowbreak{}iverses} \\
(Q4) & \texttt{Q\_one\_pa\allowbreak{}rticle\_f\allowbreak{}lat\_disp\allowbreak{}ersion}, \texttt{Q\_one\_pa\allowbreak{}rticle\_m\allowbreak{}aps}, \texttt{Q\_one\_pa\allowbreak{}rticle\_K\allowbreak{}rein\_sig\allowbreak{}natures}, \texttt{Q\_one\_pa\allowbreak{}rticle\_g\allowbreak{}eneral\_f\allowbreak{}ield} & \texttt{Q.one\_pa\allowbreak{}rticle\_m\allowbreak{}aps}, \texttt{Q.one\_pa\allowbreak{}rticle\_K\allowbreak{}rein\_ine\allowbreak{}rtia}, \texttt{compare.\allowbreak{}theory.o\allowbreak{}ne\_parti\allowbreak{}cle} \\
(Q5) & \texttt{Q\_Krein\_\allowbreak{}metric\_o\allowbreak{}f\_images} & \texttt{Q.T2\_ima\allowbreak{}ge\_keeps\allowbreak{}\_B} \\
\bottomrule
\end{longtable}
\endgroup

Relation to the good sector of the theory record. With $\Psi^\dagger$ realised as the Hilbert adjoint, the canonical anticommutator forces an indefinite (Krein) inner product (record \texttt{no\_posit\allowbreak{}ive\_inne\allowbreak{}r\_produc\allowbreak{}t}). The theory record constructs, in the good sector without extra-time momentum, a positive Fock representation in which $\Psi^\dagger$ is realised as $\chi B$ with $\chi$ the Hilbert adjoint (records \texttt{good\_sec\allowbreak{}tor\_posi\allowbreak{}tive\_foc\allowbreak{}k\_realis\allowbreak{}ation} and \texttt{Fock\_spa\allowbreak{}ce\_good\_\allowbreak{}sector\_e\allowbreak{}xample}). The statements Q1 to Q5 are statements about the canonical anticommutator; they neither use nor establish a positive-norm Fock space for either universe.

\section{Corollary C1: a T1 pair as the source of the field equations for \texorpdfstring{$a_4$}{a\_4}}

\textbf{Hypotheses.} (i) The gravitational field is the author's metric with the Einstein-Lovelock equations $\sum_{k=1}^3\alpha_kE_{(k)}{}^\mu{}_\nu + \Lambda\delta^\mu_\nu = \kappa T^\mu{}_\nu$ of \texttt{Revision\allowbreak{}/SPEC.md} section 5. (ii) The complete source is the sum of the classical energy-momentum tensors of the two members of a T1 pair at the same points of the same patch (dirac16complex00, or dirac16complex as a Grassmann-algebra-valued classical bilinear). (iii) Nothing else sources the metric.

\textbf{Statement.} The total source vanishes identically, so $a_4$ must solve the vacuum equations. (a) In Einstein gravity ($\alpha_1 = 1$, $\alpha_2 = \alpha_3 = 0$) there is no real solution for $H > 0$ and any $\Lambda$: the $x_4$ and $x_8$ components give $36H^2 + 2\Lambda = 0$ and $6(a_4')^2 + 6H^2 = 0$. (b) In Einstein-Lovelock gravity the linear member $a_4 = AHx_4 + a_0$ needs $V = 0$, with

\[
V = \alpha_1 - 40\alpha_2H^2 - 8A^2\alpha_2H^2 + 360\alpha_3H^4 + 144A^2\alpha_3H^4 + 72A^4\alpha_3H^4 ;
\]

in Einstein-Gauss-Bonnet gravity ($\alpha_1 = 1$, $\alpha_3 = 0$) this gives $A^2 = (1 - 40\alpha_2H^2)/(8\alpha_2H^2)$, real if and only if $0 < \alpha_2H^2 \leq 1/40$, with $\Lambda = -\sum_k\alpha_kE_{(k)}{}^{x_4}{}_{x_4}$.

\textbf{Proof.} (T1d) gives $T + T' = 0$, so the right-hand side of the field equations is zero, and (a) and (b) are the vacuum statements of \texttt{Revision\allowbreak{}/field\_e\allowbreak{}quations\allowbreak{}\_a4} (records below). QED.

\begingroup
\small
\begin{longtable}{@{}>{\raggedright\arraybackslash}p{0.267\linewidth}>{\raggedright\arraybackslash}p{0.267\linewidth}>{\raggedright\arraybackslash}p{0.267\linewidth}@{}}
\toprule
\textbf{statement} & \textbf{Wolfram records} & \textbf{sympy records} \\
\midrule
\endfirsthead
\toprule
\textbf{statement} & \textbf{Wolfram records} & \textbf{sympy records} \\
\midrule
\endhead
(a) & \texttt{einstein\allowbreak{}\_no\_vacu\allowbreak{}um\_solut\allowbreak{}ion} & \texttt{einstein\allowbreak{}\_no\_vacu\allowbreak{}um} \\
(b) & \texttt{linear\_m\allowbreak{}ember\_va\allowbreak{}cuum\_fac\allowbreak{}tor}, \texttt{einstein\allowbreak{}\_gauss\_b\allowbreak{}onnet\_va\allowbreak{}cuum\_lin\allowbreak{}ear} & \texttt{linear\_m\allowbreak{}ember\_va\allowbreak{}cuum\_fac\allowbreak{}tor}, \texttt{einstein\allowbreak{}\_gauss\_b\allowbreak{}onnet\_va\allowbreak{}cuum\_lin\allowbreak{}ear} \\
\bottomrule
\end{longtable}
\endgroup

Consequence (exact, under (i) to (iii)): a T1 pair cannot by itself be the source of the author's metric in Einstein gravity. This is a statement about sources in one common geometry; it is not a derivation that such a geometry, or such a pair, is created. For a T2 pair the sources do not cancel: the mirror partner carries the pulled-back energy-momentum tensor $R_8TR_8$ of the $(-m,\lambda)$ configuration. For two independently quantised universes C1 does not apply (Q3).

\section{The Kohn-Sham level T3 (to be added)}

\texttt{Revision\allowbreak{}/SPEC.md} section 9 asks for T3: at the Kohn-Sham level, the block map with the transformed boundary conditions maps the instantaneous Kohn-Sham problem with $(m,\lambda)$ to the one with $(-m,+\lambda)$ with equal energies. Status: T3 is NOT established in this document. Its exact proof is owned by \texttt{Revision\allowbreak{}/pairing\allowbreak{}/kohn\_sh\allowbreak{}am/} (a later step of the Revision work), which does not exist at the time of writing; T3 is to be added there and to this document.

The Kohn-Sham results of \texttt{Revision\allowbreak{}/kohn\_sh\allowbreak{}am/} exist, and the following ingredients of T3 are already verified there by exact checks of the Wolfram and sympy Kohn-Sham theory reports (all PASS; the boundary construction is labelled ASSUMED by that record):

\begin{itemize}
\item The block form of $\Gamma$: $\sigma_2h_j(M,k)\sigma_2 = h_{-j}(-M,k)$. It maps a solution of block type $j$ with mass $M$ onto block type $-j$ with mass $-M$ at the same $k$ and $\varepsilon$, and it exchanges the two brane parities $\chi_2(0) = 0$ and $\chi_1(0) = 0$; $\Gamma$ maps the block $(j,s_2,s_3)$ onto $(-j,s_2,s_3)$ with the $2 \times 2$ entry $s_2\sigma_2$.
\item The mirror reflection $P_A: \chi(y) \mapsto \sigma_3\chi(-y) = \gamma^8\chi(-y)$ maps $(M(y), v(y), \kappa(y))$ to $(-M(-y), v(-y), \kappa(-y))$. Under $P_A$ the densities $n$ and $Q$ are even and $S$ and the $y$-current odd, so the doubled problem of the ASSUMED Z2 orbifold is $P_A$-symmetric if and only if the mirror copy carries $(-m,+\lambda)$: $M_{\mathrm{eff}} = m + \frac{15}{16}\lambda S$ is odd and $v_v = -\lambda n/16$ even.
\item The ASSUMED Z2 mirror gives $(1 \mp \gamma^8)\chi(0) = 0$, i.e. $\chi_2(0) = 0$ (even parity) or $\chi_1(0) = 0$ (odd parity); both kill the $y$-current.
\end{itemize}

A numerical self-test of the Rust solver, labelled in its report as NOT a proof of T3: $(m, \lambda, \text{tip } \theta = 0)$ and $(-m, \lambda, \text{tip } \theta = \pi)$, solved independently, give the same sorted levels, occupations, Kohn-Sham energies and energy-momentum integrals and opposite $S$ (worst deviation 9.95e-14, tolerance 1e-9). For $N = 8$, $\lambda = 0.01946$ and $a_{4,0} = 1$: $E_{KS}$ = -9.868426190876e-4 versus -9.868426190868e-4 (80 levels), while the negative control with the untransformed tip $b(-L) = 0$ gives $E_{KS}$ = -2.0145243719e0. For $N = 136$ and $\lambda = 0.0009298$: $E_{KS}$ = 3.239294915318e1 in both runs (160 levels), negative control 2.9283751318e1. The units are those of \texttt{Revision\allowbreak{}/kohn\_sh\allowbreak{}am/ks-th\allowbreak{}eory.jso\allowbreak{}n}.

\begingroup
\small
\begin{longtable}{@{}>{\raggedright\arraybackslash}p{0.267\linewidth}>{\raggedright\arraybackslash}p{0.267\linewidth}>{\raggedright\arraybackslash}p{0.267\linewidth}@{}}
\toprule
\textbf{item} & \textbf{Wolfram records} & \textbf{sympy or Rust records} \\
\midrule
\endfirsthead
\toprule
\textbf{item} & \textbf{Wolfram records} & \textbf{sympy or Rust records} \\
\midrule
\endhead
block form of $\Gamma$ & \texttt{block\_Ga\allowbreak{}mma\_map}, \texttt{blocks\_r\allowbreak{}elation\_\allowbreak{}to\_Gamma} & \texttt{block\_Ga\allowbreak{}mma\_map}, \texttt{blocks\_r\allowbreak{}elation\_\allowbreak{}to\_Gamma} \\
mirror reflection and parities & \texttt{bc\_mirro\allowbreak{}r\_map\_PA}, \texttt{bc\_mirro\allowbreak{}r\_pariti\allowbreak{}es\_of\_de\allowbreak{}nsities} & \texttt{bc\_mirro\allowbreak{}r\_map\_PA}, \texttt{bc\_mirro\allowbreak{}r\_pariti\allowbreak{}es\_of\_de\allowbreak{}nsities} \\
brane conditions (ASSUMED) & \texttt{bc\_brane\allowbreak{}\_parity\_\allowbreak{}conditio\allowbreak{}ns} & \texttt{bc\_brane\allowbreak{}\_parity\_\allowbreak{}conditio\allowbreak{}ns} \\
numerical self-test (not a proof) & none & \texttt{t3\_block\allowbreak{}\_map\_sol\allowbreak{}ver\_self\allowbreak{}test} \\
\bottomrule
\end{longtable}
\endgroup

Interpretation (labelled, not a theorem): in its parameters the Kohn-Sham-level map is of the T2 type, $(m,\lambda) \to (-m,+\lambda)$ at equal energies, not of the T1 type.

\section{Verification records}

\subsection{Reports and their counts}

Every check has a name, a verdict and a detail. Counts at the time of writing (the publication test re-reads them):

\begingroup
\small
\begin{longtable}{@{}>{\raggedright\arraybackslash}p{0.200\linewidth}>{\raggedright\arraybackslash}p{0.200\linewidth}>{\raggedright\arraybackslash}p{0.200\linewidth}>{\raggedright\arraybackslash}p{0.200\linewidth}@{}}
\toprule
\textbf{report} & \textbf{checks} & \textbf{PASS} & \textbf{FAIL} \\
\midrule
\endfirsthead
\toprule
\textbf{report} & \textbf{checks} & \textbf{PASS} & \textbf{FAIL} \\
\midrule
\endhead
\texttt{Revision\allowbreak{}/pairing\allowbreak{}/reports\allowbreak{}/wolfram\allowbreak{}-pairing\allowbreak{}.json} & 99 & 99 & 0 \\
\texttt{Revision\allowbreak{}/pairing\allowbreak{}/reports\allowbreak{}/python-\allowbreak{}pairing.\allowbreak{}json} & 64 & 64 & 0 \\
\texttt{Revision\allowbreak{}/algebra\allowbreak{}/reports\allowbreak{}/wolfram\allowbreak{}-algebra\allowbreak{}.json} & 45 & 45 & 0 \\
\texttt{Revision\allowbreak{}/algebra\allowbreak{}/reports\allowbreak{}/python-\allowbreak{}algebra.\allowbreak{}json} & 35 & 35 & 0 \\
\texttt{Revision\allowbreak{}/theory/\allowbreak{}reports/\allowbreak{}wolfram-\allowbreak{}field-th\allowbreak{}eory.jso\allowbreak{}n} & 84 & 84 & 0 \\
\texttt{Revision\allowbreak{}/theory/\allowbreak{}reports/\allowbreak{}python-f\allowbreak{}ield-the\allowbreak{}ory.json} & 70 & 70 & 0 \\
\texttt{Revision\allowbreak{}/field\_e\allowbreak{}quations\allowbreak{}\_a4/repo\allowbreak{}rts/wolf\allowbreak{}ram-a4-r\allowbreak{}eport.js\allowbreak{}on} & 47 & 47 & 0 \\
\texttt{Revision\allowbreak{}/field\_e\allowbreak{}quations\allowbreak{}\_a4/repo\allowbreak{}rts/pyth\allowbreak{}on-a4-re\allowbreak{}port.jso\allowbreak{}n} & 61 & 61 & 0 \\
\texttt{Revision\allowbreak{}/kohn\_sh\allowbreak{}am/repor\allowbreak{}ts/ks-th\allowbreak{}eory-wol\allowbreak{}fram.jso\allowbreak{}n} & 46 & 46 & 0 \\
\texttt{Revision\allowbreak{}/kohn\_sh\allowbreak{}am/repor\allowbreak{}ts/ks-th\allowbreak{}eory-pyt\allowbreak{}hon.json} & 57 & 57 & 0 \\
\texttt{Revision\allowbreak{}/kohn\_sh\allowbreak{}am/repor\allowbreak{}ts/ks-ru\allowbreak{}st-solve\allowbreak{}r.json} & 40 & 40 & 0 \\
\bottomrule
\end{longtable}
\endgroup

The pairing theorems are verified in four gravitational fields: the author's metric on the patch $z \in (0,\pi/2)$ (primordial); the same metric on the mirror patch; a general diagonal (4,4) field $e^a{}_\mu = h_a(x_1,\dots,x_8)\delta^a_\mu$ with eight arbitrary functions of all eight coordinates (diagonal8); and a general non-diagonal field at one point (pointwise), with exact rational $e^a{}_\mu$ and $\partial_\rho e^a{}_\mu$, $\det e = 1411/6561$, a non-diagonal metric and 448 of the 512 connection components nonzero; plus the linearity (kernel) argument that covers every field. The canonical connection of each test field satisfies the vielbein postulate exactly. The Wolfram side and the sympy side share no code; the sympy side re-builds the gammas from the author's tau formulas and compares them entry by entry with \texttt{Revision\allowbreak{}/algebra\allowbreak{}/gammas.\allowbreak{}json}.

\subsection{Checks per theorem}

The tables below list all 99 checks of the Wolfram pairing report and all 64 checks of the sympy pairing report (section 9.1), grouped by theorem; every one has the verdict PASS. Among the common sympy checks, the three comparison checks for T1, T2 and Q record that the sympy side confirms each Wolfram theorem independently (with 17, 13 and 8 named sympy checks), and the comparison check for the limits records that the two independently written lists of what the theorems do not establish cover the same topics; section 11.2 merges them.

\textbf{T1, Wolfram} (51 checks):

\begingroup
\small
\begin{longtable}{@{}>{\raggedright\arraybackslash}p{0.267\linewidth}>{\raggedright\arraybackslash}p{0.267\linewidth}>{\raggedright\arraybackslash}p{0.267\linewidth}@{}}
\toprule
\textbf{check} & \textbf{check} & \textbf{check} \\
\midrule
\endfirsthead
\toprule
\textbf{check} & \textbf{check} & \textbf{check} \\
\midrule
\endhead
\texttt{T1\_kerne\allowbreak{}l\_scalar} & \texttt{T1\_kerne\allowbreak{}l\_kineti\allowbreak{}c} & \texttt{T1\_kerne\allowbreak{}l\_connec\allowbreak{}tion} \\
\texttt{T1\_kerne\allowbreak{}l\_field\_\allowbreak{}equation} & \texttt{T1\_linea\allowbreak{}rity\_arg\allowbreak{}ument} & \texttt{T1\_scala\allowbreak{}r\_and\_ki\allowbreak{}netic\_pr\allowbreak{}imordial\allowbreak{}\_commuti\allowbreak{}ng} \\
\texttt{T1\_Lagra\allowbreak{}ngian\_pr\allowbreak{}imordial\allowbreak{}\_commuti\allowbreak{}ng} & \texttt{T1\_Lagra\allowbreak{}ngian\_ne\allowbreak{}gative\_c\allowbreak{}ontrols\_\allowbreak{}primordi\allowbreak{}al\_commu\allowbreak{}ting} & \texttt{T1\_field\allowbreak{}\_equatio\allowbreak{}n\_covari\allowbreak{}ance\_pri\allowbreak{}mordial\_\allowbreak{}commutin\allowbreak{}g} \\
\texttt{T1\_Euler\allowbreak{}\_Lagrang\allowbreak{}e\_derive\allowbreak{}d\_primor\allowbreak{}dial\_com\allowbreak{}muting} & \texttt{T1\_Euler\allowbreak{}\_Lagrang\allowbreak{}e\_map\_pr\allowbreak{}imordial\allowbreak{}\_commuti\allowbreak{}ng} & \texttt{T1\_energ\allowbreak{}y\_moment\allowbreak{}um\_primo\allowbreak{}rdial\_co\allowbreak{}mmuting} \\
\texttt{T1\_curre\allowbreak{}nt\_primo\allowbreak{}rdial\_co\allowbreak{}mmuting} & \texttt{T1\_gener\allowbreak{}al\_poten\allowbreak{}tial\_pri\allowbreak{}mordial\_\allowbreak{}commutin\allowbreak{}g} & \texttt{T1\_scala\allowbreak{}r\_and\_ki\allowbreak{}netic\_pr\allowbreak{}imordial\allowbreak{}\_grassma\allowbreak{}nn} \\
\texttt{T1\_Lagra\allowbreak{}ngian\_pr\allowbreak{}imordial\allowbreak{}\_grassma\allowbreak{}nn} & \texttt{T1\_Lagra\allowbreak{}ngian\_ne\allowbreak{}gative\_c\allowbreak{}ontrols\_\allowbreak{}primordi\allowbreak{}al\_grass\allowbreak{}mann} & \texttt{T1\_field\allowbreak{}\_equatio\allowbreak{}n\_covari\allowbreak{}ance\_pri\allowbreak{}mordial\_\allowbreak{}grassman\allowbreak{}n} \\
\texttt{T1\_Euler\allowbreak{}\_Lagrang\allowbreak{}e\_derive\allowbreak{}d\_primor\allowbreak{}dial\_gra\allowbreak{}ssmann} & \texttt{T1\_Euler\allowbreak{}\_Lagrang\allowbreak{}e\_map\_pr\allowbreak{}imordial\allowbreak{}\_grassma\allowbreak{}nn} & \texttt{T1\_energ\allowbreak{}y\_moment\allowbreak{}um\_primo\allowbreak{}rdial\_gr\allowbreak{}assmann} \\
\texttt{T1\_curre\allowbreak{}nt\_primo\allowbreak{}rdial\_gr\allowbreak{}assmann} & \texttt{T1\_scala\allowbreak{}r\_and\_ki\allowbreak{}netic\_di\allowbreak{}agonal8\_\allowbreak{}commutin\allowbreak{}g} & \texttt{T1\_Lagra\allowbreak{}ngian\_di\allowbreak{}agonal8\_\allowbreak{}commutin\allowbreak{}g} \\
\texttt{T1\_Lagra\allowbreak{}ngian\_ne\allowbreak{}gative\_c\allowbreak{}ontrols\_\allowbreak{}diagonal\allowbreak{}8\_commut\allowbreak{}ing} & \texttt{T1\_field\allowbreak{}\_equatio\allowbreak{}n\_covari\allowbreak{}ance\_dia\allowbreak{}gonal8\_c\allowbreak{}ommuting} & \texttt{T1\_Euler\allowbreak{}\_Lagrang\allowbreak{}e\_derive\allowbreak{}d\_diagon\allowbreak{}al8\_comm\allowbreak{}uting} \\
\texttt{T1\_Euler\allowbreak{}\_Lagrang\allowbreak{}e\_map\_di\allowbreak{}agonal8\_\allowbreak{}commutin\allowbreak{}g} & \texttt{T1\_energ\allowbreak{}y\_moment\allowbreak{}um\_diago\allowbreak{}nal8\_com\allowbreak{}muting} & \texttt{T1\_curre\allowbreak{}nt\_diago\allowbreak{}nal8\_com\allowbreak{}muting} \\
\texttt{T1\_scala\allowbreak{}r\_and\_ki\allowbreak{}netic\_di\allowbreak{}agonal8\_\allowbreak{}grassman\allowbreak{}n} & \texttt{T1\_Lagra\allowbreak{}ngian\_di\allowbreak{}agonal8\_\allowbreak{}grassman\allowbreak{}n} & \texttt{T1\_Lagra\allowbreak{}ngian\_ne\allowbreak{}gative\_c\allowbreak{}ontrols\_\allowbreak{}diagonal\allowbreak{}8\_grassm\allowbreak{}ann} \\
\texttt{T1\_field\allowbreak{}\_equatio\allowbreak{}n\_covari\allowbreak{}ance\_dia\allowbreak{}gonal8\_g\allowbreak{}rassmann} & \texttt{T1\_Euler\allowbreak{}\_Lagrang\allowbreak{}e\_derive\allowbreak{}d\_diagon\allowbreak{}al8\_gras\allowbreak{}smann} & \texttt{T1\_Euler\allowbreak{}\_Lagrang\allowbreak{}e\_map\_di\allowbreak{}agonal8\_\allowbreak{}grassman\allowbreak{}n} \\
\texttt{T1\_energ\allowbreak{}y\_moment\allowbreak{}um\_diago\allowbreak{}nal8\_gra\allowbreak{}ssmann} & \texttt{T1\_curre\allowbreak{}nt\_diago\allowbreak{}nal8\_gra\allowbreak{}ssmann} & \texttt{T1\_scala\allowbreak{}r\_and\_ki\allowbreak{}netic\_po\allowbreak{}intwise\_\allowbreak{}commutin\allowbreak{}g} \\
\texttt{T1\_Lagra\allowbreak{}ngian\_po\allowbreak{}intwise\_\allowbreak{}commutin\allowbreak{}g} & \texttt{T1\_Lagra\allowbreak{}ngian\_ne\allowbreak{}gative\_c\allowbreak{}ontrols\_\allowbreak{}pointwis\allowbreak{}e\_commut\allowbreak{}ing} & \texttt{T1\_field\allowbreak{}\_equatio\allowbreak{}n\_covari\allowbreak{}ance\_poi\allowbreak{}ntwise\_c\allowbreak{}ommuting} \\
\texttt{T1\_energ\allowbreak{}y\_moment\allowbreak{}um\_point\allowbreak{}wise\_com\allowbreak{}muting} & \texttt{T1\_curre\allowbreak{}nt\_point\allowbreak{}wise\_com\allowbreak{}muting} & \texttt{T1\_scala\allowbreak{}r\_and\_ki\allowbreak{}netic\_po\allowbreak{}intwise\_\allowbreak{}grassman\allowbreak{}n} \\
\texttt{T1\_Lagra\allowbreak{}ngian\_po\allowbreak{}intwise\_\allowbreak{}grassman\allowbreak{}n} & \texttt{T1\_Lagra\allowbreak{}ngian\_ne\allowbreak{}gative\_c\allowbreak{}ontrols\_\allowbreak{}pointwis\allowbreak{}e\_grassm\allowbreak{}ann} & \texttt{T1\_field\allowbreak{}\_equatio\allowbreak{}n\_covari\allowbreak{}ance\_poi\allowbreak{}ntwise\_g\allowbreak{}rassmann} \\
\texttt{T1\_energ\allowbreak{}y\_moment\allowbreak{}um\_point\allowbreak{}wise\_gra\allowbreak{}ssmann} & \texttt{T1\_curre\allowbreak{}nt\_point\allowbreak{}wise\_gra\allowbreak{}ssmann} & \texttt{Gamma\_pr\allowbreak{}operties} \\
\bottomrule
\end{longtable}
\endgroup

\textbf{T1, sympy} (23 checks):

\begingroup
\small
\begin{longtable}{@{}>{\raggedright\arraybackslash}p{0.267\linewidth}>{\raggedright\arraybackslash}p{0.267\linewidth}>{\raggedright\arraybackslash}p{0.267\linewidth}@{}}
\toprule
\textbf{check} & \textbf{check} & \textbf{check} \\
\midrule
\endfirsthead
\toprule
\textbf{check} & \textbf{check} & \textbf{check} \\
\midrule
\endhead
\texttt{T1.gener\allowbreak{}al\_field\allowbreak{}.matrix\_\allowbreak{}identiti\allowbreak{}es} & \texttt{T1.gener\allowbreak{}al\_field\allowbreak{}.random\_\allowbreak{}instance\allowbreak{}.commuti\allowbreak{}ng} & \texttt{T1.gener\allowbreak{}al\_field\allowbreak{}.random\_\allowbreak{}instance\allowbreak{}.grassma\allowbreak{}nn} \\
\texttt{T1.metri\allowbreak{}c.commut\allowbreak{}ing.S\_in\allowbreak{}variant} & \texttt{T1.metri\allowbreak{}c.commut\allowbreak{}ing.lagr\allowbreak{}angian} & \texttt{T1.metri\allowbreak{}c.commut\allowbreak{}ing.nega\allowbreak{}tive\_con\allowbreak{}trols} \\
\texttt{T1.metri\allowbreak{}c.commut\allowbreak{}ing.eule\allowbreak{}r\_lagran\allowbreak{}ge\_deriv\allowbreak{}ed} & \texttt{T1.metri\allowbreak{}c.commut\allowbreak{}ing.eule\allowbreak{}r\_lagran\allowbreak{}ge\_map} & \texttt{T1.metri\allowbreak{}c.commut\allowbreak{}ing.emt} \\
\texttt{T1.metri\allowbreak{}c.commut\allowbreak{}ing.pair\allowbreak{}\_total\_e\allowbreak{}mt\_zero} & \texttt{T1.metri\allowbreak{}c.commut\allowbreak{}ing.curr\allowbreak{}ent} & \texttt{T1.metri\allowbreak{}c.grassm\allowbreak{}ann.S\_in\allowbreak{}variant} \\
\texttt{T1.metri\allowbreak{}c.grassm\allowbreak{}ann.lagr\allowbreak{}angian} & \texttt{T1.metri\allowbreak{}c.grassm\allowbreak{}ann.nega\allowbreak{}tive\_con\allowbreak{}trols} & \texttt{T1.metri\allowbreak{}c.grassm\allowbreak{}ann.eule\allowbreak{}r\_lagran\allowbreak{}ge\_deriv\allowbreak{}ed} \\
\texttt{T1.metri\allowbreak{}c.grassm\allowbreak{}ann.eule\allowbreak{}r\_lagran\allowbreak{}ge\_map} & \texttt{T1.metri\allowbreak{}c.grassm\allowbreak{}ann.emt} & \texttt{T1.metri\allowbreak{}c.grassm\allowbreak{}ann.pair\allowbreak{}\_total\_e\allowbreak{}mt\_zero} \\
\texttt{T1.metri\allowbreak{}c.grassm\allowbreak{}ann.curr\allowbreak{}ent} & \texttt{T1.metri\allowbreak{}c.commut\allowbreak{}ing.gene\allowbreak{}ral\_pote\allowbreak{}ntial} & \texttt{T1.metri\allowbreak{}c.grassm\allowbreak{}ann.gene\allowbreak{}ral\_pote\allowbreak{}ntial} \\
\texttt{T1.diago\allowbreak{}nal8.com\allowbreak{}muting} & \texttt{T1.diago\allowbreak{}nal8.gra\allowbreak{}ssmann} &  \\
\bottomrule
\end{longtable}
\endgroup

\textbf{T2, Wolfram} (28 checks):

\begingroup
\small
\begin{longtable}{@{}>{\raggedright\arraybackslash}p{0.267\linewidth}>{\raggedright\arraybackslash}p{0.267\linewidth}>{\raggedright\arraybackslash}p{0.267\linewidth}@{}}
\toprule
\textbf{check} & \textbf{check} & \textbf{check} \\
\midrule
\endfirsthead
\toprule
\textbf{check} & \textbf{check} & \textbf{check} \\
\midrule
\endhead
\texttt{T2\_Pn\_in\allowbreak{}\_Pin44} & \texttt{T2\_Pn\_co\allowbreak{}vers\_the\allowbreak{}\_reflect\allowbreak{}ion} & \texttt{T2\_chara\allowbreak{}cter\_of\_\allowbreak{}Pn} \\
\texttt{T2\_Gamma\allowbreak{}\_times\_P\allowbreak{}n\_is\_gam\allowbreak{}ma\_n} & \texttt{T2\_kerne\allowbreak{}ls\_gamma\allowbreak{}\_n\_with\_\allowbreak{}frame\_re\allowbreak{}flection} & \texttt{T2\_frame\allowbreak{}\_reflect\allowbreak{}ion\_poin\allowbreak{}twise\_x1\allowbreak{}\_commuti\allowbreak{}ng} \\
\texttt{T2\_frame\allowbreak{}\_reflect\allowbreak{}ion\_poin\allowbreak{}twise\_x1\allowbreak{}\_grassma\allowbreak{}nn} & \texttt{T2\_frame\allowbreak{}\_reflect\allowbreak{}ion\_poin\allowbreak{}twise\_x2\allowbreak{}\_commuti\allowbreak{}ng} & \texttt{T2\_frame\allowbreak{}\_reflect\allowbreak{}ion\_poin\allowbreak{}twise\_x2\allowbreak{}\_grassma\allowbreak{}nn} \\
\texttt{T2\_frame\allowbreak{}\_reflect\allowbreak{}ion\_poin\allowbreak{}twise\_x3\allowbreak{}\_commuti\allowbreak{}ng} & \texttt{T2\_frame\allowbreak{}\_reflect\allowbreak{}ion\_poin\allowbreak{}twise\_x3\allowbreak{}\_grassma\allowbreak{}nn} & \texttt{T2\_frame\allowbreak{}\_reflect\allowbreak{}ion\_poin\allowbreak{}twise\_x4\allowbreak{}\_commuti\allowbreak{}ng} \\
\texttt{T2\_frame\allowbreak{}\_reflect\allowbreak{}ion\_poin\allowbreak{}twise\_x4\allowbreak{}\_grassma\allowbreak{}nn} & \texttt{T2\_frame\allowbreak{}\_reflect\allowbreak{}ion\_poin\allowbreak{}twise\_x5\allowbreak{}\_commuti\allowbreak{}ng} & \texttt{T2\_frame\allowbreak{}\_reflect\allowbreak{}ion\_poin\allowbreak{}twise\_x5\allowbreak{}\_grassma\allowbreak{}nn} \\
\texttt{T2\_frame\allowbreak{}\_reflect\allowbreak{}ion\_poin\allowbreak{}twise\_x6\allowbreak{}\_commuti\allowbreak{}ng} & \texttt{T2\_frame\allowbreak{}\_reflect\allowbreak{}ion\_poin\allowbreak{}twise\_x6\allowbreak{}\_grassma\allowbreak{}nn} & \texttt{T2\_frame\allowbreak{}\_reflect\allowbreak{}ion\_poin\allowbreak{}twise\_x7\allowbreak{}\_commuti\allowbreak{}ng} \\
\texttt{T2\_frame\allowbreak{}\_reflect\allowbreak{}ion\_poin\allowbreak{}twise\_x7\allowbreak{}\_grassma\allowbreak{}nn} & \texttt{T2\_frame\allowbreak{}\_reflect\allowbreak{}ion\_poin\allowbreak{}twise\_x8\allowbreak{}\_commuti\allowbreak{}ng} & \texttt{T2\_frame\allowbreak{}\_reflect\allowbreak{}ion\_poin\allowbreak{}twise\_x8\allowbreak{}\_grassma\allowbreak{}nn} \\
\texttt{T2\_mirro\allowbreak{}r\_is\_iso\allowbreak{}metry} & \texttt{T2\_mirro\allowbreak{}r\_Lagran\allowbreak{}gian\_com\allowbreak{}muting} & \texttt{T2\_mirro\allowbreak{}r\_energy\allowbreak{}\_momentu\allowbreak{}m\_and\_cu\allowbreak{}rrent\_co\allowbreak{}mmuting} \\
\texttt{T2\_mirro\allowbreak{}r\_Euler\_\allowbreak{}Lagrange\allowbreak{}\_commuti\allowbreak{}ng} & \texttt{T2\_mirro\allowbreak{}r\_Lagran\allowbreak{}gian\_gra\allowbreak{}ssmann} & \texttt{T2\_mirro\allowbreak{}r\_energy\allowbreak{}\_momentu\allowbreak{}m\_and\_cu\allowbreak{}rrent\_gr\allowbreak{}assmann} \\
\texttt{T2\_mirro\allowbreak{}r\_Euler\_\allowbreak{}Lagrange\allowbreak{}\_grassma\allowbreak{}nn} &  &  \\
\bottomrule
\end{longtable}
\endgroup

\textbf{T2, sympy} (16 checks):

\begingroup
\small
\begin{longtable}{@{}>{\raggedright\arraybackslash}p{0.267\linewidth}>{\raggedright\arraybackslash}p{0.267\linewidth}>{\raggedright\arraybackslash}p{0.267\linewidth}@{}}
\toprule
\textbf{check} & \textbf{check} & \textbf{check} \\
\midrule
\endfirsthead
\toprule
\textbf{check} & \textbf{check} & \textbf{check} \\
\midrule
\endhead
\texttt{geometry\allowbreak{}.brane\_d\allowbreak{}egenerat\allowbreak{}e} & \texttt{geometry\allowbreak{}.mirror\_\allowbreak{}isometry} & \texttt{T2.gener\allowbreak{}al\_field\allowbreak{}.matrix\_\allowbreak{}identiti\allowbreak{}es} \\
\texttt{T2.gener\allowbreak{}al\_field\allowbreak{}.reflect\allowbreak{}ion\_tabl\allowbreak{}e} & \texttt{T2.metri\allowbreak{}c.commut\allowbreak{}ing.refl\allowbreak{}ection\_t\allowbreak{}able} & \texttt{T2.metri\allowbreak{}c.commut\allowbreak{}ing.eule\allowbreak{}r\_lagran\allowbreak{}ge\_map} \\
\texttt{T2.metri\allowbreak{}c.commut\allowbreak{}ing.emt} & \texttt{T2.metri\allowbreak{}c.commut\allowbreak{}ing.curr\allowbreak{}ent} & \texttt{T2.metri\allowbreak{}c.commut\allowbreak{}ing.S\_od\allowbreak{}d} \\
\texttt{T2.metri\allowbreak{}c.grassm\allowbreak{}ann.refl\allowbreak{}ection\_t\allowbreak{}able} & \texttt{T2.metri\allowbreak{}c.grassm\allowbreak{}ann.eule\allowbreak{}r\_lagran\allowbreak{}ge\_map} & \texttt{T2.metri\allowbreak{}c.grassm\allowbreak{}ann.emt} \\
\texttt{T2.metri\allowbreak{}c.grassm\allowbreak{}ann.curr\allowbreak{}ent} & \texttt{T2.metri\allowbreak{}c.grassm\allowbreak{}ann.S\_od\allowbreak{}d} & \texttt{T2.diago\allowbreak{}nal8.fra\allowbreak{}me\_refle\allowbreak{}ctions.c\allowbreak{}ommuting} \\
\texttt{T2.diago\allowbreak{}nal8.fra\allowbreak{}me\_refle\allowbreak{}ctions.g\allowbreak{}rassmann} &  &  \\
\bottomrule
\end{longtable}
\endgroup

\textbf{Q, Wolfram} (10 checks):

\begingroup
\small
\begin{longtable}{@{}>{\raggedright\arraybackslash}p{0.267\linewidth}>{\raggedright\arraybackslash}p{0.267\linewidth}>{\raggedright\arraybackslash}p{0.267\linewidth}@{}}
\toprule
\textbf{check} & \textbf{check} & \textbf{check} \\
\midrule
\endfirsthead
\toprule
\textbf{check} & \textbf{check} & \textbf{check} \\
\midrule
\endhead
\texttt{Q\_Krein\_\allowbreak{}metric\_o\allowbreak{}f\_images} & \texttt{Q\_symple\allowbreak{}ctic\_ker\allowbreak{}nel\_comm\allowbreak{}uting} & \texttt{Q\_genera\allowbreak{}tors\_of\_\allowbreak{}the\_imag\allowbreak{}e\_commut\allowbreak{}ing} \\
\texttt{Q\_symple\allowbreak{}ctic\_ker\allowbreak{}nel\_gras\allowbreak{}smann} & \texttt{Q\_genera\allowbreak{}tors\_of\_\allowbreak{}the\_imag\allowbreak{}e\_grassm\allowbreak{}ann} & \texttt{Q\_one\_pa\allowbreak{}rticle\_f\allowbreak{}lat\_disp\allowbreak{}ersion} \\
\texttt{Q\_one\_pa\allowbreak{}rticle\_m\allowbreak{}aps} & \texttt{Q\_one\_pa\allowbreak{}rticle\_K\allowbreak{}rein\_sig\allowbreak{}natures} & \texttt{Q\_one\_pa\allowbreak{}rticle\_g\allowbreak{}eneral\_f\allowbreak{}ield} \\
\texttt{Q\_no\_ide\allowbreak{}ntificat\allowbreak{}ion\_of\_i\allowbreak{}ndepende\allowbreak{}nt\_unive\allowbreak{}rses} &  &  \\
\bottomrule
\end{longtable}
\endgroup

\textbf{Q, sympy} (8 checks):

\begingroup
\small
\begin{longtable}{@{}>{\raggedright\arraybackslash}p{0.267\linewidth}>{\raggedright\arraybackslash}p{0.267\linewidth}>{\raggedright\arraybackslash}p{0.267\linewidth}@{}}
\toprule
\textbf{check} & \textbf{check} & \textbf{check} \\
\midrule
\endfirsthead
\toprule
\textbf{check} & \textbf{check} & \textbf{check} \\
\midrule
\endhead
\texttt{Q.canoni\allowbreak{}cal\_anti\allowbreak{}commutat\allowbreak{}or} & \texttt{Q.image\_\allowbreak{}krein\_me\allowbreak{}tric} & \texttt{Q.image\_\allowbreak{}own\_quan\allowbreak{}tisation} \\
\texttt{Q.image\_\allowbreak{}generato\allowbreak{}rs\_same\_\allowbreak{}dynamics} & \texttt{Q.no\_can\allowbreak{}cellatio\allowbreak{}n\_indepe\allowbreak{}ndent\_un\allowbreak{}iverses} & \texttt{Q.one\_pa\allowbreak{}rticle\_m\allowbreak{}aps} \\
\texttt{Q.one\_pa\allowbreak{}rticle\_K\allowbreak{}rein\_ine\allowbreak{}rtia} & \texttt{Q.T2\_ima\allowbreak{}ge\_keeps\allowbreak{}\_B} &  \\
\bottomrule
\end{longtable}
\endgroup

\textbf{Common checks (gammas, geometry, comparison with the Wolfram record), sympy} (17 checks):

\begingroup
\small
\begin{longtable}{@{}>{\raggedright\arraybackslash}p{0.267\linewidth}>{\raggedright\arraybackslash}p{0.267\linewidth}>{\raggedright\arraybackslash}p{0.267\linewidth}@{}}
\toprule
\textbf{check} & \textbf{check} & \textbf{check} \\
\midrule
\endfirsthead
\toprule
\textbf{check} & \textbf{check} & \textbf{check} \\
\midrule
\endhead
\texttt{gammas.c\allowbreak{}lifford} & \texttt{\detokenize{gammas.C}} & \texttt{gammas.G\allowbreak{}amma} \\
\texttt{gammas.G\allowbreak{}amma\_ant\allowbreak{}icommute\allowbreak{}s} & \texttt{\detokenize{gammas.B}} & \texttt{gammas.e\allowbreak{}qual\_wol\allowbreak{}fram\_fix\allowbreak{}ture} \\
\texttt{geometry\allowbreak{}.omega\_a\allowbreak{}ntisymme\allowbreak{}tric} & \texttt{geometry\allowbreak{}.gamma\_c\allowbreak{}ovariant\allowbreak{}ly\_const\allowbreak{}ant} & \texttt{geometry\allowbreak{}.spin\_co\allowbreak{}nnection\allowbreak{}\_compone\allowbreak{}nts} \\
\texttt{compare.\allowbreak{}theory.s\allowbreak{}tatus} & \texttt{compare.\allowbreak{}theory.r\allowbreak{}eflectio\allowbreak{}n\_table} & \texttt{compare.\allowbreak{}theory.k\allowbreak{}rein\_sig\allowbreak{}ns} \\
\texttt{compare.\allowbreak{}theory.o\allowbreak{}ne\_parti\allowbreak{}cle} & \texttt{compare.\allowbreak{}theory.t\allowbreak{}heorem\_T\allowbreak{}1} & \texttt{compare.\allowbreak{}theory.t\allowbreak{}heorem\_T\allowbreak{}2} \\
\texttt{compare.\allowbreak{}theory.t\allowbreak{}heorem\_Q} & \texttt{compare.\allowbreak{}theory.n\allowbreak{}ot\_estab\allowbreak{}lished} &  \\
\bottomrule
\end{longtable}
\endgroup

\textbf{Common checks (fixture, test fields, connections, determinism), Wolfram} (10 checks):

\begingroup
\small
\begin{longtable}{@{}>{\raggedright\arraybackslash}p{0.267\linewidth}>{\raggedright\arraybackslash}p{0.267\linewidth}>{\raggedright\arraybackslash}p{0.267\linewidth}@{}}
\toprule
\textbf{check} & \textbf{check} & \textbf{check} \\
\midrule
\endfirsthead
\toprule
\textbf{check} & \textbf{check} & \textbf{check} \\
\midrule
\endhead
\texttt{fixture\_\allowbreak{}Clifford\allowbreak{}\_relatio\allowbreak{}n} & \texttt{fixture\_\allowbreak{}definiti\allowbreak{}ons} & \texttt{C\_and\_B\_\allowbreak{}basic} \\
\texttt{primordi\allowbreak{}al\_vielb\allowbreak{}ein} & \texttt{pointwis\allowbreak{}e\_field\_\allowbreak{}is\_gener\allowbreak{}al} & \texttt{connecti\allowbreak{}on\_primo\allowbreak{}rdial} \\
\texttt{connecti\allowbreak{}on\_mirro\allowbreak{}r\_patch} & \texttt{connecti\allowbreak{}on\_diago\allowbreak{}nal8} & \texttt{connecti\allowbreak{}on\_point\allowbreak{}wise} \\
\texttt{theory\_f\allowbreak{}ile\_dete\allowbreak{}rministi\allowbreak{}c} &  &  \\
\bottomrule
\end{longtable}
\endgroup

\subsection{Re-verification for this document}

On 2026-10-01 both pairing verifiers were re-run, for this document, on a copy of the working tree in a scratch directory: \texttt{verify\_p\allowbreak{}airing.w\allowbreak{}ls} gave 99 of 99 PASS in 80.5 s and \texttt{check\_pa\allowbreak{}iring.py} gave 64 of 64 PASS in 143.6 s (the two ran concurrently on the development machine); the regenerated \texttt{wolfram-\allowbreak{}pairing.\allowbreak{}json}, \texttt{pairing-\allowbreak{}theory.j\allowbreak{}son} and \texttt{python-p\allowbreak{}airing.j\allowbreak{}son} were byte-identical to the committed files.

\section{Reproduction}

From the repository root (the sympy pairing checker reads the Wolfram theory record for its final comparison, so the Wolfram verifier runs first; the PDF command is one command, the backslash is the line continuation of a POSIX shell):

\begin{Verbatim}[fontsize=\small]
wolframscript -file Revision/algebra/wolfram/verify_algebra.wls
python Revision/algebra/python/check_algebra.py
wolframscript -file Revision/theory/wolfram/verify_field_theory.wls
python Revision/theory/python/check_field_theory.py
wolframscript -file Revision/pairing/wolfram/verify_pairing.wls
python Revision/pairing/python/check_pairing.py
wolframscript -file Revision/field_equations_a4/wolfram/verify_field_equations_a4.wls
python Revision/field_equations_a4/python/check_field_equations_a4.py
wolframscript -file Revision/kohn_sham/theory/verify_ks_theory.wls
python Revision/kohn_sham/theory/check_ks_theory.py
python scripts/build_provenance_pdf.py Revision/docs/PAIR_CREATION_PROOFS.md \
    --developer-layout --specifications Revision/pdf-specifications.json
python -m unittest Revision/tests/test_pair_creation_proofs_publication.py -v
\end{Verbatim}

The PDF build runs the Markdown-to-LaTeX builder twice (with different hash seeds), runs pdflatex three times into each of two fresh directories, requires warning-free logs and byte-identical PDFs, and checks the PDF against its entry \texttt{pair-cre\allowbreak{}ation-pr\allowbreak{}oofs} in \texttt{Revision\allowbreak{}/pdf-spe\allowbreak{}cificati\allowbreak{}ons.json} (the Revision registry; the registry of the earlier stages is not touched).

\section{What is proved and what is not}

\subsection{Proved}

Exactly, for both fields (dirac16complex with Grassmann components, dirac16complex00 with commuting components) unless stated, under the hypotheses stated with each result:

\begin{enumerate}
\item T1 (section 4): $\mathcal{L}_{m,\lambda}[\Gamma\Psi] = -\mathcal{L}_{-m,-\lambda}[\Psi]$; the solutions of the $(m,\lambda)$ theory and of the $(-m,-\lambda)$ theory correspond one to one; $T \to -T$ and $J \to -J$; the pair has zero total energy-momentum, current and charge as classical bilinears; in every gravitational field taken as a fixed background.
\item T2 (section 5): $\Gamma$ combined with a Pin(4,4) reflection of character $-1$ gives $\mathcal{L}_{m,\lambda}[\gamma^n\Psi; R_ne] = +\mathcal{L}_{-m,\lambda}[\Psi; e]$; in the author's field the mirror across the Z2 brane $z = \pi/2$ (ASSUMED construction) pairs a $(-m,\lambda)$ solution on the patch with an $(m,\lambda)$ solution on the mirror patch at EQUAL energy-momentum and charge.
\item Q (section 6, dirac16complex only): the chirality image carries the Krein metric $-B$ and is the same quantum system re-labelled; an independently quantised $(-m,-\lambda)$ universe carries $+B$ and cannot be identified with it; the generators of two independent universes add without cancelling; the $+m$ and $-m$ one-particle spectra are identical; the T2 image keeps $+B$.
\item C1 (section 7): a T1 pair as the complete classical source of the author's metric is a zero source; in Einstein gravity the author's metric then has no solution for $H > 0$.
\end{enumerate}

\subsection{Not established}

\begin{enumerate}
\item No creation process: the theorems map solutions to solutions and quantities to quantities; nothing in these equations produces a universe, a pair of universes or a change of the number of universes, and no transition between ``no universe'' and ``two universes'', no initial state, no vacuum decay and no tunnelling process is derived.
\item No rate, probability or amplitude: no transition amplitude, probability, cross-section, rate or Bogoliubov coefficient for creating universes of masses $\{+m, -m\}$ is computed or implied; no wave function of the universe and no path integral is part of these theorems.
\item No dynamical necessity: no equation and no conservation law forces the partner to exist; a single universe with mass $+m$ is an equally valid solution without its partner. T1 and T2 are correspondences between solutions of two parameter sets, not a mechanism.
\item T1 is not a symmetry of a single theory: it changes the parameters and the sign of the action, and for $\lambda \neq 0$ it pairs $(m,\lambda)$ with $(-m,-\lambda)$, not $+m$ with $-m$ at fixed coupling. The pure pairing $(m,\lambda) \to (-m,\lambda)$ is T2, at equal, not opposite, energy-momentum: T2 does not reverse the sign of $T$.
\item The vanishing total energy-momentum and charge of a T1 pair holds for classical bilinears and as an operator identity within ONE quantum system; it does not hold for two independently quantised universes, whose generators add without cancelling.
\item Test-field statement: the gravitational field is fixed and the same for both members. The back-reaction through the field equations for $a_4$ is not part of the theorems; corollary C1 is the only statement about it, and it concerns the sum of classical sources in one common geometry.
\item The Z2 brane: the mirror across $z = \pi/2$ uses the ASSUMED Z2 construction; the metric is degenerate there ($g_{88} = 0$ and $\sqrt{|g|} = 0$), and no junction condition, brane tension or matching of the field across the brane is derived.
\item Quantum positivity: in flat 4+4 space every energy eigenspace of the one-particle Hamiltonian computed exactly at the samples of section 6.3 has Krein inertia (4,4); a positive-norm Fock space for either universe is not established by these theorems (the good-sector construction of the theory record is a separate result, not used here).
\item dirac16complex00 is a classical field: no quantum statement is made for it.
\item The Kohn-Sham level T3 is not established in this document (section 8).
\end{enumerate}

\subsection{The answer to the request}

The request ``PROVE that Universes of masses \{+mass, -mass\} are created in pairs'' is answered exactly as follows, for dirac16complex and for dirac16complex00. Proved: universes of masses $+m$ and $-m$ come in pairs in the precise sense of the theorems T1 and T2. To every solution with mass $m$ the explicit maps $\Gamma$ (with $\lambda \to -\lambda$, opposite energy-momentum and charge) and $\gamma^8$ with the Z2 mirror (same $\lambda$, equal energy-momentum) assign a solution with mass $-m$, in the author's primordial gravitational field and, for T1, in every gravitational field; at the quantum level of dirac16complex the chirality partner is the same quantum system with the Krein metric $-B$, and two independent universes of masses $+m$ and $-m$ have identical spectra and no cancelling generators. Not proved: that such universes are CREATED, in pairs or otherwise. No creation process, rate or amplitude follows from these equations.

\end{document}
